\chapter{自由电子传播与有限时间跃迁}\label{chap:electron-finite-time}

电子与环境的态空间、谱和耦合已经确定，但矩阵元还没有告诉我们一段有限作用时间后的概率。
电子的自由色散会改变波包宽度，相互作用则让振幅在不同环境占据数之间流动。
将这些变化放进同一个联合传播子，才能区分相干交换、有限时间谱宽与长时间的恒定跃迁率。

\section{连续通道核与自由传播}

含时 Schr\"odinger 方程的解由联合传播子表示。将它在环境量子数基底中展开，可以追踪不同交换路径；对连续末态求和与长时间极限则决定何时出现 Fermi 黄金律。

有限时间实验首先由联合传播子的连续通道核决定。取电子正能连续基 $|\vb{p},s\rangle$，并在辅助系统谱空间中取完备广义基 $|\mu\rangle_X$。以下以 $\int\dd\mu_X(\mu)$ 统一表示离散求和与连续谱积分，并以 $\delta_X(\mu,\nu)$ 表示相对于该谱测度的恒等核，即
\[
\int\dd\mu_X(\mu)\,\delta_X(\mu,\nu)f(\mu)=f(\nu).
\]
令 $T=t_f-t_i$，定义有限时间联合通道核
\begin{equation}
{
\mathscr U_{\mu\nu}^{s_f s_i}
(\vb{p}_f,\vb{p}_i;T)
\equiv
\langle \vb{p}_f,s_f;\mu|
\hat U(t_f,t_i)
|\vb{p}_i,s_i;\nu\rangle .}
\label{eq:exact_finite_time_channel_kernel}
\end{equation}
$\mathscr U$ 是分布意义下的振幅密度：电子反冲、自旋翻转、任意辅助谱结构以及有限作用时间都仍保留在它的矩阵元中。联合幺正性给出连续通道正交归一恒等式
\begin{equation}
{
\sum_{s_f}\int\dd^3p_f\int\dd\mu_X(\mu)\,
\bigl[\mathscr U_{\mu\nu}^{s_f s_i}
(\vb{p}_f,\vb{p}_i;T)\bigr]^*
\mathscr U_{\mu\nu'}^{s_f s_i'}
(\vb{p}_f,\vb{p}_i';T)
=
\delta_X(\nu,\nu')\delta_{s_i s_i'}
\delta^{(3)}(\vb{p}_i-\vb{p}_i').}
\label{eq:exact_channel_kernel_unitarity}
\end{equation}
这条恒等式给出了任何模式截断、边带模型或有效主方程都必须保持的概率约束。

更一般地，把初态联合密度算符的连续核定义为
\[
\rho_{eX,i}^{s_i s_i';\nu\nu'}(\vb{p}_i,\vb{p}_i')
\equiv
\langle\vb{p}_i,s_i;\nu|\hat\rho_{eX}(t_i)
|\vb{p}_i',s_i';\nu'\rangle .
\]
对辅助末态取偏迹后，电子输出密度核为
\begin{equation}
{
\begin{aligned}
\rho_{e,f}^{s_f s_f'}(\vb{p}_f,\vb{p}_f')
={}&\sum_{s_i,s_i'}
\int\dd^3p_i\dd^3p_i'
\int\dd\mu_X(\mu)\dd\mu_X(\nu)\dd\mu_X(\nu')\,
\mathscr U_{\mu\nu}^{s_f s_i}
(\vb{p}_f,\vb{p}_i;T)\\
&\times
\rho_{eX,i}^{s_i s_i';\nu\nu'}(\vb{p}_i,\vb{p}_i')
\bigl[\mathscr U_{\mu\nu'}^{s_f' s_i'}
(\vb{p}_f',\vb{p}_i';T)\bigr]^* .
\end{aligned}}
\label{eq:exact_reduced_electron_channel_kernel}
\end{equation}
因此电子输入相干、辅助系统相干、初始电子--环境相关、反冲、自旋以及所有未观测辅助末态都由同一个积分核传播。只有当
\[
\hat\rho_{eX}(t_i)=\hat\rho_e(t_i)\otimes\hat\rho_X(t_i)
\]
时，初态核才因子化为
\[
\rho_{eX,i}^{s_i s_i';\nu\nu'}
=\rho_\ee^{s_i s_i'}\rho_X^{\nu\nu'},
\]
固定同一个 $\hat\rho_X(t_i)$ 后，所得电子演化是\emph{完全正保迹}（CPTP）映射。保迹指归一化态演化后总概率仍为一；完全正指即使电子预先与另一个不参与相互作用的参考系统纠缠，只对电子施加此映射，联合态仍保持非负。它们来自联合幺正演化和环境偏迹。若初始电子--环境相关存在，仍使用\eqnrefs{eq:exact_reduced_electron_channel_kernel} 直接传播联合初态，而不把约化动力学强行写成与初始相关无关的量子信道。

实验建模通常还会截断模式、Fock 数或电子边带。令 $\hat P_{\rm keep}$ 投影到准备保留的联合末态子空间，$\hat Q_{\rm keep}=\mathbb I-\hat P_{\rm keep}$；令 $\hat\Pi_{\rm in}$ 投影到实际允许的联合输入工作窗。定义有限时间最坏情形泄漏量
\begin{equation}
{
\epsilon_{\rm proj}(T)
\equiv
\bigl\|\hat Q_{\rm keep}\hat U(t_f,t_i)\hat\Pi_{\rm in}\bigr\|^2.}
\label{eq:exact_projection_leakage_bound}
\end{equation}
对任何支撑在 $\hat\Pi_{\rm in}$ 内的归一化输入态，所有被删通道的总概率都不超过 $\epsilon_{\rm proj}(T)$。这里截取的是完整传播子产生的末态。若实验对 $P_{\rm keep}$ 后选择，重新归一化给条件态，成功概率也必须保留。若要用压缩哈密顿量 $P_{\rm keep}HP_{\rm keep}$ 代替完整动力学，还需检验被删通道的虚跃迁和累计相位；末态泄漏本身不控制这些效应。这个判据不依赖被保留通道被称为腔模、极化激元、PINEM 边带还是 Bragg 格点。

\begin{theorem}[联合通道核的幺正性、偏迹与泄漏]
设 $U$ 是联合 Hilbert 空间的幺正算符，初态为归一化正迹类算符。通道核满足\eqnrefs{eq:exact_channel_kernel_unitarity}，约化核为\eqnrefs{eq:exact_reduced_electron_channel_kernel}。对入射窗，最坏末态泄漏概率为 $\|Q_{\rm keep}U\Pi_{\rm in}\|^2$。固定环境且初态因子化时，约化动力学为 CPTP 信道；预先相关初态不能直接定义任意输入态上的同一信道。
\end{theorem}
\begin{proof}
辅助系统的广义完备关系写成
\[
\mathbb I_X=\int\dd\mu_X(\mu)|\mu\rangle\langle\mu|,
\]
电子正能一粒子扇区的完备关系为
\[
\mathbb I_e=\sum_s\int\dd^3p\,|\vb{p},s\rangle\langle\vb{p},s|.
\]
因此联合完备关系是
\[
\mathbb I_{eX}
=\sum_{s_f}\int\dd^3p_f\int\dd\mu_X(\mu)
|\vb{p}_f,s_f;\mu\rangle
\langle\vb{p}_f,s_f;\mu|.
\]
在两个任意入射广义基矢之间插入
$\hat U^\dagger\hat U=\mathbb I_{eX}$：
\begin{align*}
&\langle\vb{p}_i,s_i;\nu|\hat U^\dagger\hat U
|\vb{p}_i',s_i';\nu'\rangle\\
&=\sum_{s_f}\int\dd^3p_f\int\dd\mu_X(\mu)\,
\langle\vb{p}_i,s_i;\nu|\hat U^\dagger
|\vb{p}_f,s_f;\mu\rangle
\langle\vb{p}_f,s_f;\mu|\hat U
|\vb{p}_i',s_i';\nu'\rangle\\
&=\sum_{s_f}\int\dd^3p_f\int\dd\mu_X(\mu)\,
\bigl[\mathscr U_{\mu\nu}^{s_f s_i}
(\vb{p}_f,\vb{p}_i;T)\bigr]^*
\mathscr U_{\mu\nu'}^{s_f s_i'}
(\vb{p}_f,\vb{p}_i';T).
\end{align*}
广义基的归一化给
\[
\langle\vb{p}_i,s_i;\nu|\vb{p}_i',s_i';\nu'\rangle
=\delta_X(\nu,\nu')\delta_{s_i s_i'}
\delta^{(3)}(\vb{p}_i-\vb{p}_i'),
\]
从而得到\eqnrefs{eq:exact_channel_kernel_unitarity}。

再从任意联合初态开始。按定义，
\begin{align*}
\hat\rho_{eX}(t_i)
={}&\sum_{s_i,s_i'}
\int\dd^3p_i\dd^3p_i'
\int\dd\mu_X(\nu)\dd\mu_X(\nu')\,
\rho_{eX,i}^{s_i s_i';\nu\nu'}(\vb{p}_i,\vb{p}_i')\\
&\times
|\vb{p}_i,s_i;\nu\rangle
\langle\vb{p}_i',s_i';\nu'|.
\end{align*}
联合末态为
$\hat\rho_{eX}(t_f)=\hat U\hat\rho_{eX}(t_i)\hat U^\dagger$。电子密度核由辅助系统偏迹给出：
\[
\rho_{e,f}^{s_f s_f'}(\vb{p}_f,\vb{p}_f')
=\int\dd\mu_X(\mu)
\langle\vb{p}_f,s_f;\mu|\hat\rho_{eX}(t_f)
|\vb{p}_f',s_f';\mu\rangle .
\]
代入联合初态展开，并分别使用
\[
\langle\vb{p}_f,s_f;\mu|\hat U
|\vb{p}_i,s_i;\nu\rangle
=\mathscr U_{\mu\nu}^{s_f s_i}(\vb{p}_f,\vb{p}_i;T)
\]
与其复共轭，逐项积分便得到\eqnrefs{eq:exact_reduced_electron_channel_kernel}。

截断误差也可直接从同一传播子得到。对任意满足
$|\Psi_i\rangle=\hat\Pi_{\rm in}|\Psi_i\rangle$ 的归一化纯态，被删末态的概率为
\begin{align*}
P_{\rm out}
&=\langle\Psi_i|\hat U^\dagger\hat Q_{\rm keep}\hat U|\Psi_i\rangle\\
&=\|\hat Q_{\rm keep}\hat U\hat\Pi_{\rm in}|\Psi_i\rangle\|^2\\
&\le
\|\hat Q_{\rm keep}\hat U\hat\Pi_{\rm in}\|^2
=\epsilon_{\rm proj}(T).
\end{align*}
混合输入态可写成纯态凸组合，因此同一上界逐项成立并在求和后保持。对纯态输入取上确界还给
\[
\sup_{\Psi_i\in\operatorname{Ran}\Pi_{\rm in}}P_{\rm out}
=\epsilon_{\rm proj}(T),
\]
所以\eqnrefs{eq:exact_projection_leakage_bound} 正是所选输入工作窗上的最坏情形总泄漏概率，而不是某个具体模式模型的经验参数。

最后对电子末态取迹：
\begin{align*}
\operatorname{Tr}_e\hat\rho_e(t_f)
&=\sum_{s_f}\int\dd^3p_f\,
\rho_{e,f}^{s_f s_f}(\vb{p}_f,\vb{p}_f)\\
&=\operatorname{Tr}_{eX}
\left(\hat U\hat\rho_{eX}(t_i)\hat U^\dagger\right)\\
&=\operatorname{Tr}_{eX}\hat\rho_{eX}(t_i)=1.
\end{align*}
若进一步令 $\hat\rho_{eX}(t_i)=\hat\rho_e\otimes\hat\rho_X$，
将固定环境态作谱分解 $\rho_X=\sum_r\lambda_r|r\rangle\langle r|$，
其中 $\lambda_r\ge0$、$\sum_r\lambda_r=1$。定义电子算符
\[
K_{\mu r}=\sqrt{\lambda_r}\,{}_X\langle\mu|U|r\rangle_X.
\]
由环境完备关系和 $U^\dagger U=I$，得到
\[
\Phi(\rho_e)=\sum_r\int\dd\mu_X(\mu)K_{\mu r}\rho_eK_{\mu r}^\dagger,
\qquad
\sum_r\int\dd\mu_X(\mu)K_{\mu r}^\dagger K_{\mu r}=I_e.
\]
第二式保证保迹。对任意参考空间 $R$ 及正联合输入，
第一式中的每项在 $K_{\mu r}\otimes I_R$ 的左右作用后仍为正，
所以该映射完全正；这就是这里所用的 Kraus 表示。
连续基时将这些式子理解为弱算符积分，或先用归一化波包离散化再取极限。
这一构造为任意电子输入态指定了同一个环境初态，因而定义了同一量子信道。
存在初始相关时，则必须另外说明联合态怎样随电子输入而制备；仅给电子边缘态不能确定输出。
\end{proof}


\begin{proposition}[末态后选择的成功概率与状态误差]
令 $\sigma=U\rho U^\dagger$ 为归一化联合末态，$P=P_{\rm keep}$ 为正交投影，
$q=\operatorname{Tr}[(I-P)\sigma]<1$。保留事件的成功概率为 $s=1-q$，
其条件态为 $\sigma_{\rm c}=P\sigma P/s$。定义迹距离
$D(\sigma,\tau)=\tfrac12\|\sigma-\tau\|_1$，则
\[
D(\sigma,\sigma_{\rm c})\le\sqrt q,\qquad
|\operatorname{Tr}(M\sigma)-\operatorname{Tr}(M\sigma_{\rm c})|
\le\sqrt q\quad(0\le M\le I).
\]
若入射态支撑在 $\Pi_{\rm in}$ 内，便可用
$q\le\epsilon_{\rm proj}(T)$ 给出统一上界。
对电子或环境再取偏迹后，迹距离界仍成立。
\end{proposition}
\begin{proof}
取 $\sigma$ 的一个归一化纯化 $|\Psi\rangle$，辅助纯化空间记为 $R$。
投影成功后的归一化纯化为
\[
|\Phi\rangle=\frac{(P\otimes I_R)|\Psi\rangle}{\sqrt s},\qquad
\langle\Psi|\Phi\rangle=\sqrt s.
\]
两个纯态之差只作用于它们张成的二维空间。
将第二个态写成 $\sqrt s|\Psi\rangle+\sqrt q|\Psi_\perp\rangle$，
可见差矩阵的迹为零、行列式为 $-q$，故其两个非零本征值为 $\pm\sqrt q$。
两纯态的迹距离于是等于 $\sqrt q$。
偏迹不会增加迹距离：对任意效应 $M$，约化后的概率差等于联合空间中
$M\otimes I_R$ 的概率差；对效应取上确界便得这一收缩性。
先对 $R$ 求偏迹，再按需对未观测系统求偏迹，即得所给状态界。
最后，迹为零的厄米差算符的正、负部分具有相同迹，
所以任意 $0\le M\le I$ 的概率差不超过其迹距离。
\end{proof}

这个界比较完整末态与对完整末态作后选择得到的条件态。
它没有把完整 $U$ 换成截断哈密顿量的传播子。
电子可以暂时离开保留空间，再带着相位返回；
出射时看不到泄漏，并不等于中途经过的通道没有作用。

\begin{exbox}[零末态泄漏仍会改变干涉相位]
取三个正交联合通道 $|0\rangle,|1\rangle,|a\rangle$，
保留 $P=|0\rangle\langle0|+|1\rangle\langle1|$，删去 $|a\rangle$。
在旋转绘景中设
\[
H=\hbar g\bigl(|1\rangle\langle a|+|a\rangle\langle1|\bigr),\qquad g>0.
\]
以 $|+\rangle=(|0\rangle+|1\rangle)/\sqrt2$ 入射，
求 $T=\pi/g$ 时的泄漏和电子干涉读出，
并与直接用 $PHP$ 演化的预测比较。

\textbf{解。} $|0\rangle$ 不参与耦合；在 $|1\rangle,|a\rangle$ 子空间中，
$H/\hbar=g\sigma_x$，所以
\[
U(t)|1\rangle=\cos(gt)|1\rangle-\ii\sin(gt)|a\rangle.
\]
任意保留输入的最坏泄漏为 $\epsilon_{\rm proj}(t)=\sin^2(gt)$。
在所给末时刻，它等于零，却有
\[
U(T)|+\rangle=\frac{|0\rangle-|1\rangle}{\sqrt2}=| -\rangle.
\]
完整模型对效应 $M_+=|+\rangle\langle+|$ 给出概率零。
另一方面，$PHP=0$，其传播子在保留空间中为恒等算符，
预测同一效应的概率为一。误差达到最大值，尽管末态泄漏完全消失。

差别来自耦合通道的一次完整往返。
这里两个保留分量之间的相对相位可由干涉读出，不能当作公共相位删去。
若要删去 $|a\rangle$，需要保留它产生的有效能移与相位，
或直接比较完整传播子与候选有效传播子；后一比较由下面的 Duhamel 恒等式控制。
\end{exbox}


弱耦合 FGR、离散反冲格、PINEM/QPINEM 边带以及开放系统主方程都从这套有限时间通道结构向下约化：先按实际谱和测量分辨率选择通道并检查\eqnrefs{eq:exact_projection_leakage_bound}，再比较有限时间失谐、矩阵元变化、反冲曲率与环境记忆尺度，最后才实施对应近似。新的材料模式、非高斯光场或不同电子色散只改变 $\mathscr U$ 的具体矩阵元与谱测度。

这里的 Hilbert 空间分解以正能一电子扇区与辅助系统为对象；若研究条件允许显著的电子--正电子对产生、粒子数改变或必须保留虚粒子多体扇区，则联合动力学应采用《高等物理》相对论量子场部分的完整哈密顿量，再把相应多粒子散射振幅投影到实验可见通道。

\begin{exbox}[窄带高斯波包与三阶相位余项]
\label{sec:feqo-dispersion-expansion}
取中心动量 $\mathbf p_0=p_0\mathbf e_z$，正能色散的完整二阶展开为
\begin{equation}
E(\mathbf p)=E_0+v_0\delta p_z+
\frac{\delta p_x^2+\delta p_y^2}{2\gamma_0m_{\rm e}}
+\frac{\delta p_z^2}{2\gamma_0^3m_{\rm e}}+R_3(\delta\mathbf p).
\label{eq:feqo-dispersion-hessian}
\end{equation}
这由 $\partial_i\partial_jE=c^2\delta_{ij}/E-c^4p_ip_j/E^3$ 在中心处取值。
固定正能自旋，把纵向动量包络取为
\begin{equation}
a(p)=(2\pi\sigma_p^2)^{-1/4}
\exp[-(p-p_0)^2/(4\sigma_p^2)].
\label{eq:feqo-gaussian-momentum}
\end{equation}
在窄带 $\sigma_p\ll p_0$ 内，负向尾概率
$\tfrac12\operatorname{erfc}(p_0/(\sqrt2\sigma_p))$ 可忽略。
令 $m_\parallel=\gamma_0^3m_{\rm e}$，
使用前面色散 Hessian 的纵向分量，并令 $\xi=z-v_0t$。
抽去共同载波后，Fourier 积分中的指数为
\[
-\frac{\delta p^2}{4\sigma_p^2}
\left(1+\ii\frac{2\sigma_p^2t}{m_\parallel\hbar}\right)
+\frac{\ii\delta p\xi}{\hbar}.
\]
定义 $\sigma_{z0}=\hbar/(2\sigma_p)$、
$\tau=\hbar t/(2m_\parallel\sigma_{z0}^2)$。
配方并完成高斯积分，得到归一化包络
\begin{align}
\psi(\xi,t)&=(2\pi\sigma_{z0}^2)^{-1/4}(1+\ii\tau)^{-1/2}
\exp\left[-\frac{\xi^2}{4\sigma_{z0}^2(1+\ii\tau)}\right],\label{eq:feqo-gaussian-prop-result}\\
|\psi(\xi,t)|^2&=
\frac{\exp[-\xi^2/(2\sigma_{z0}^2(1+\tau^2))]}
{\sqrt{2\pi}\sigma_{z0}\sqrt{1+\tau^2}},\label{eq:feqo-gaussian-density}\\
\sigma_z(t)^2&=\sigma_{z0}^2+
\frac{\hbar^2t^2}{4m_\parallel^2\sigma_{z0}^2}.\label{eq:feqo-gaussian-spreading}
\end{align}
$\tau=1$ 时方差增至两倍，标准差增至 $\sqrt2$ 倍。
对 $\SI{200}{\kilo\electronvolt}$ 电子，$\sigma_{z0}=\SI{1}{\nano\metre}$
给此时间 $t_*=\SI{4.65e-14}{\second}$、漂移长度 $v_0t_*=\SI{9.70}{\micro\metre}$；
初宽增至 $\SI{10}{\nano\metre}$ 时，这两个尺度都增大一百倍。

高斯没有有限支集。实际误差必须同时计入主体内的相位余项和主体外的尾部。
取 $|\delta p|\le r\sigma_p$，令 $M_3$ 是此窗内 $|E'''|$ 的上界。
尾概率为 $\eta_r=\operatorname{erfc}(r/\sqrt2)$，则
\begin{equation}
\|(U_{\rm ex}-U_2)\psi_0\|
\le\frac{tM_3(r\sigma_p)^3}{6\hbar}+2\sqrt{\eta_r}.
\label{eq:feqo-gaussian-propagation-tail-control}
\end{equation}
主体界由 $|\ee^{\ii x}-1|\le|x|$，尾界由两个幺正算符的范数和给出。
若旋量随动量变化，还需独立加入冻结旋量的幅度误差
$\sup_{|\delta p|\le r\sigma_p}\|u_s(p)-u_s(p_0)\|$。
因此窄带模型的适用范围由累计谱相位、尾概率及旋量变化共同确定。
\end{exbox}

\begin{figure}[htbp]
\centering
\includegraphics[width=0.76\linewidth]{relativistic_effective_mass.pdf}
\caption{同一相对论色散的横向与纵向曲率质量比，分别为 $\gamma$ 和 $\gamma^3$。}
\label{fig:relativistic-effective-mass}
\end{figure}
\begin{figure}[htbp]
\centering
\includegraphics[width=0.76\linewidth]{dispersion_control.pdf}
\caption{200\,\si{\kilo\electronvolt} 电子的精确色散与局部一阶、二阶展开。高阶残差还需乘传播时间除以 $\hbar$，才能判断传播相位误差。}
\label{fig:dispersion-control-modern}
\end{figure}


\subsection{传播子与平均场}

求解时常要略去哈密顿量中较小的项。不过，一个瞬时很小的项经过长时间积累，也可能明显改变相位。判断近似是否可靠，需要比较两种哈密顿量从同一个初态产生的演化。为此设
$\hat H(t)=\hat H_0(t)+\hat R(t)$，其中 $\hat H_0$ 是待检验的参考动力学，$\hat R$ 是被省略的余项；两者从同一初时刻 $t_i$ 出发产生传播子 $\hat U$ 与 $\hat U_0$。它们的差精确满足
\begin{equation}
{
\hat U(t_f,t_i)-\hat U_0(t_f,t_i)
=-\frac{\ii}{\hbar}\int_{t_i}^{t_f}\dd t\,
\hat U(t_f,t)\hat R(t)\hat U_0(t,t_i).}
\label{eq:duhamel_master_identity}
\end{equation}
因此，对实际准备的初态子空间投影 $\hat\Pi_{\rm in}$，幺正性给出
\begin{equation}
{
\|[\hat U-\hat U_0]\hat\Pi_{\rm in}\|
\le
\frac1\hbar\int_{t_i}^{t_f}\dd t\,
\|\hat R(t)\hat U_0(t,t_i)\hat\Pi_{\rm in}\|.}
\label{eq:duhamel_projected_bound}
\end{equation}
这里的范数衡量两种输出态之间可能出现的最大差别，$\hat\Pi_{\rm in}$ 只允许取实际制备范围内的初态。右边的 $\hat U_0(t,t_i)$ 先把初态传播到时刻 $t$，$\hat R(t)$ 再作用于这一时刻的态，最后对作用时间积分。因而，一个被舍弃的项是否可忽略，同时取决于它有多大、作用了多久，以及电子是否进入它能显著作用的态。

\begin{theorem}[Duhamel 恒等式与入射窗误差]
设 $H_0(t)$ 自伴并生成幺正传播子，$R(t)=R(t)^\dagger$ 有界且算符范数可积，$H=H_0+R$。则\eqnrefs{eq:duhamel_master_identity}成立，任意入射投影满足\eqnrefs{eq:duhamel_projected_bound}。若归一化态的传播误差至多 $\epsilon$，任意效应 $0\le E\le I$ 的概率误差至多 $2\epsilon$。无界余项需验证共同域及被积范数可积，单有有限初态投影不充分。
\end{theorem}
\begin{proof}
先在定理的有界余项条件下计算。若 $H_0$ 无界，以下求导可在共同演化域上理解，再由积分恒等式延拓；有限入射投影本身不能替代定义域条件。令完整与参考哈密顿量分别为

\begin{equation*}
H(t)=H_0(t)+R(t),
\end{equation*}
对应传播子满足
\begin{align*}
\ii\hbar\partial_tU(t,t_i)&=H(t)U(t,t_i),\\
\ii\hbar\partial_tU_0(t,t_i)&=H_0(t)U_0(t,t_i).
\end{align*}
关键是对中间时刻 $t$ 的混合乘积
$F(t)=U(t_f,t)U_0(t,t_i)$ 求导。由于传播子的第一、第二时间参数导数符号相反，
\begin{align*}
\partial_tU(t_f,t)&=+\frac{\ii}{\hbar}U(t_f,t)H(t),\\
\partial_tU_0(t,t_i)&=-\frac{\ii}{\hbar}H_0(t)U_0(t,t_i),
\end{align*}
于是
\begin{align*}
\frac{\dd F}{\dd t}
&=\frac{\ii}{\hbar}U(t_f,t)
[H(t)-H_0(t)]U_0(t,t_i)\\
&=\frac{\ii}{\hbar}U(t_f,t)R(t)U_0(t,t_i).
\end{align*}
从 $t_i$ 积分到 $t_f$：
\begin{align*}
F(t_f)-F(t_i)
&=U_0(t_f,t_i)-U(t_f,t_i)\\
&=\frac{\ii}{\hbar}
\int_{t_i}^{t_f}\dd t\,
U(t_f,t)R(t)U_0(t,t_i).
\end{align*}
移项即得 Duhamel 母恒等式。若只关心一个实际入射工作窗 $\Pi_{\rm in}$，则右乘投影：
\begin{equation*}
[U-U_0]\Pi_{\rm in}
=-\frac{\ii}{\hbar}
\int\dd t\,U(t_f,t)R(t)U_0(t,t_i)\Pi_{\rm in}.
\end{equation*}
对算符范数使用三角不等式与次乘性，若完整传播子幺正，$\|U(t_f,t)\|=1$，便有
\begin{equation*}
\|[U-U_0]\Pi_{\rm in}\|
\le\frac1\hbar\int\dd t\,
\|R(t)U_0(t,t_i)\Pi_{\rm in}\|.
\end{equation*}
该表达式保留了参考动力学怎样把初始工作窗拖入别的 Hilbert 空间扇区的信息，比粗略的
$\int\|R(t)\|\dd t/\hbar$ 更紧。若右端小于 $\epsilon$，则对所有
$|\psi\rangle\in\operatorname{Ran}\Pi_{\rm in}$、$\|\psi\|=1$ 都有
\begin{equation*}
\|(U-U_0)|\psi\rangle\|\le\epsilon,
\end{equation*}
因此任何归一化末态投影概率的差都至多为 $O(\epsilon)$；这给出有限时间近似控制量的共同来源。

最后一项的明确常数来自 $\langle u|E|u\rangle-\langle v|E|v\rangle=\langle u-v|E|u\rangle+\langle v|E|u-v\rangle$，其中 $\|u\|=\|v\|=1$。两项各不超过 $\|u-v\|$。
\end{proof}


现在把这个工具用于量子场到经典场的过渡。若第 $\lambda$ 个玻色模式围绕相干幅 $\alpha_\lambda(t)=\langle\hat a_\lambda\rangle$，则将模式算符分解为平均值与零均值涨落
\[
\hat a_\lambda=\alpha_\lambda+\delta\hat a_\lambda,
\qquad
\langle\delta\hat a_\lambda\rangle=0.
\]
$\alpha_\lambda$ 无量纲，决定平均场；$\delta\hat a_\lambda$ 决定平均场之外仍能与电子纠缠的量子涨落。代入\eqnrefs{eq:universal_free_electron_boson_hamiltonian} 后，电子看到的平均驱动为
\begin{equation}
{
\hat V_{\rm cl}(t)
=\hbar\sum_\lambda
\left[
\alpha_\lambda(t)\hat{\mathcal G}_\lambda(t)
+\alpha_\lambda^*(t)\hat{\mathcal G}_\lambda^\dagger(t)
\right].}
\label{eq:multimode_classical_drive_operator}
\end{equation}
而剩余量子部分为
\[
\hat R_{\rm q}(t)
=\hbar\sum_\lambda
\left(
\hat{\mathcal G}_\lambda\delta\hat a_\lambda
+\hat{\mathcal G}_\lambda^\dagger\delta\hat a_\lambda^\dagger
\right).
\]
因此“光子很多”本身不是经典极限；真正需要同时保持平均驱动有限并让涨落传播子修正消失。

由于 $\hat a_\lambda$ 在完整 Fock 空间上无界，误差必须在实验实际访问的有限工作窗上定义。令 $\hat\Pi_{\Lambda,\vb{N}}$ 限制电子到能量窗 $\Lambda$，并限制各模式的涨落占据数到 $0\le n_\lambda\le N_\lambda$；再定义
\[
G_{\lambda,\Lambda}(t)
\equiv
\max\!\left\{\|\hat{\mathcal G}_\lambda(t)\hat\Pi_\Lambda\|,
\|\hat{\mathcal G}_\lambda^\dagger(t)\hat\Pi_\Lambda\|\right\}.
\]
经典极限对应联合缩放
\begin{equation}
{
|\alpha_\lambda|\to\infty,
\qquad
G_{\lambda,\Lambda}\to0,
\qquad
|\alpha_\lambda|G_{\lambda,\Lambda}=O(1),}
\label{eq:multimode_coherent_classical_scaling}
\end{equation}
即单量子场幅变小、占据数变大，而电子看到的平均驱动保持有限。若 $\hat U_{\rm full}$ 保留涨落，$\hat U_{\rm cl}$ 只保留平均场，则
\begin{equation}
{
\left\|[\hat U_{\rm full}-\hat U_{\rm cl}]\hat\Pi_{\rm in}\right\|
\le
\frac1\hbar\int_{t_i}^{t_f}\dd t\,
\|\hat R_{\rm q}(t)\hat U_{\rm cl}(t,t_i)\hat\Pi_{\rm in}\|.}
\label{eq:multimode_classical_limit_duhamel_bound}
\end{equation}
只要参考动力学在作用时间内没有显著离开所选工作窗，右端相对于平均驱动按 $\max_\lambda\sqrt{N_\lambda+1}/|\alpha_\lambda|$ 衰减；若工作窗存在尾概率，还需另加相应的尾部误差。于是电子的经典场哈密顿量为
\begin{equation}
{
\hat H_{\rm e,cl}(t)=\hat H_e+\hat V_{\rm cl}(t),}
\label{eq:general_classical_field_limit_electron_hamiltonian}
\end{equation}
但它的有限时间演化仍是完整的时间排序传播子，而不是自动等价于一阶微扰。

\begin{proposition}[多模相干平均场的有限时间误差]
对有限个线性耦合玻色模，位移初始相干态后，将哈密顿量精确分成平均场与涨落余项。令 $G_{\lambda,\Lambda}$ 同时控制 $\mathcal G_\lambda\Pi_\Lambda$ 与 $\mathcal G_\lambda^\dagger\Pi_\Lambda$ 的范数；在所用传播子和余项积分存在时，\eqnrefs{eq:multimode_classical_limit_duhamel_bound}控制平均场传播误差，窗外余项必须另行计入。
\end{proposition}
\begin{proof}
从一般线性电子--玻色相互作用
\[
\hat V(t)
=\hbar\sum_\lambda
\left[
\hat{\mathcal G}_\lambda(t)\hat a_\lambda
+\hat{\mathcal G}_\lambda^\dagger(t)\hat a_\lambda^\dagger
\right]
\]
出发。设环境初态在所保留模式上为多模相干态
\[
|\vb{\alpha}\rangle
=\bigotimes_\lambda|\alpha_\lambda\rangle,
\qquad
\hat a_\lambda|\vb{\alpha}\rangle
=\alpha_\lambda|\vb{\alpha}\rangle .
\]
与其把 $\hat a_\lambda$ 直接“替换成一个复数”，更稳妥的做法是先用多模位移算符
\[
\hat D(\vb{\alpha})
=\exp\!\left[
\sum_\lambda
(\alpha_\lambda\hat a_\lambda^\dagger-\alpha_\lambda^*\hat a_\lambda)
\right]
\]
把相干态平移到真空。Baker--Campbell--Hausdorff 公式给
\[
\hat D^\dagger\hat a_\lambda\hat D
=\hat a_\lambda+\alpha_\lambda,
\qquad
\hat D^\dagger\hat a_\lambda^\dagger\hat D
=\hat a_\lambda^\dagger+\alpha_\lambda^* .
\]
因此在位移绘景中，
\begin{align*}
\hat D^\dagger\hat V(t)\hat D
&=
\hbar\sum_\lambda
\left[
\alpha_\lambda\hat{\mathcal G}_\lambda(t)
+\alpha_\lambda^*\hat{\mathcal G}_\lambda^\dagger(t)
\right]\\
&\quad+
\hbar\sum_\lambda
\left[
\hat{\mathcal G}_\lambda(t)\hat a_\lambda
+\hat{\mathcal G}_\lambda^\dagger(t)\hat a_\lambda^\dagger
\right]\\
&\equiv
\hat V_{\rm cl}(t)+\hat R_{\rm q}(t).
\end{align*}
第一项就是正文\eqnrefs{eq:multimode_classical_drive_operator}。这一分解不是近似，而只是把平均相干幅与围绕它的量子涨落放在不同项中；近似真正发生在后面舍弃 $\hat R_{\rm q}$ 时。

为了给这个舍弃一个有限时间控制，取电子工作窗投影 $\hat\Pi_\Lambda$，并把每个保留模式的涨落 Fock 空间截到 $0\le n_\lambda\le N_\lambda$，记联合投影为
\[
\hat\Pi_{\Lambda,\vb{N}}
=\hat\Pi_\Lambda\otimes
\bigotimes_\lambda
\sum_{n=0}^{N_\lambda}|n\rangle\langle n|.
\]
定义
\[
G_{\lambda,\Lambda}(t)
\equiv
\max\!\left\{\|\hat{\mathcal G}_\lambda(t)\hat\Pi_\Lambda\|,\|\hat{\mathcal G}_\lambda^\dagger(t)\hat\Pi_\Lambda\|\right\}.
\]
由
\[
\|\hat a_\lambda\hat\Pi_{\vb{N}}\|
\le\sqrt{N_\lambda},
\qquad
\|\hat a_\lambda^\dagger\hat\Pi_{\vb{N}}\|
\le\sqrt{N_\lambda+1}
\]
得到
\[
\frac1\hbar
\|\hat R_{\rm q}(t)\hat\Pi_{\Lambda,\vb{N}}\|
\le
\sum_\lambda
G_{\lambda,\Lambda}(t)
\bigl(\sqrt{N_\lambda}+\sqrt{N_\lambda+1}\bigr).
\]
经典极限不是简单地令 $|\alpha_\lambda|\to\infty$；若同时保持微观耦合不变，平均驱动也会发散。受控缩放应保持
\[
\mathcal E_{\lambda,\rm cl}(t)
\equiv
|\alpha_\lambda|G_{\lambda,\Lambda}(t)
\]
有限，即让单量子耦合随 $|\alpha_\lambda|^{-1}$ 缩放。于是
\[
G_{\lambda,\Lambda}
=\frac{\mathcal E_{\lambda,\rm cl}}{|\alpha_\lambda|},
\]
从而
\[
\frac1\hbar
\|\hat R_{\rm q}\hat\Pi_{\Lambda,\vb{N}}\|
\le
\max_\lambda
\frac{\sqrt{N_\lambda+1}}{|\alpha_\lambda|}
\,
2\sum_\lambda\mathcal E_{\lambda,\rm cl}
\]
到一个与 $|\alpha_\lambda|$ 无关的常数因子。

令 $\hat U$ 是完整位移绘景传播子，$\hat U_{\rm cl}$ 是只保留 $\hat V_{\rm cl}$ 的传播子。Duhamel 恒等式给
\[
\hat U(t_f,t_i)-\hat U_{\rm cl}(t_f,t_i)
=
-\frac{\ii}{\hbar}
\int_{t_i}^{t_f}\dd t\,
\hat U(t_f,t)\hat R_{\rm q}(t)\hat U_{\rm cl}(t,t_i).
\]
令 $\hat Q_{\Lambda,\vb{N}}=\mathbb I-\hat\Pi_{\Lambda,\vb{N}}$，并把参考动力学离开受控工作窗后对涨落余项的贡献定义为
\[
\epsilon_{\rm cl,tail}
\equiv
\frac1\hbar\int_{t_i}^{t_f}\dd t\,
\|\hat R_{\rm q}(t)\hat Q_{\Lambda,\vb{N}}
\hat U_{\rm cl}(t,t_i)\hat\Pi_{\rm in}\|.
\]
在 Duhamel 被积函数中插入
$\hat\Pi_{\Lambda,\vb{N}}+\hat Q_{\Lambda,\vb{N}}=\mathbb I$，得到
\begin{align*}
\|(\hat U-\hat U_{\rm cl})\hat\Pi_{\rm in}\|
&\le
\frac1\hbar
\int_{t_i}^{t_f}\dd t\,
\|\hat R_{\rm q}(t)\hat\Pi_{\Lambda,\vb{N}}\|
+\epsilon_{\rm cl,tail}\\
&=
O\!\left(
\max_\lambda
\frac{\sqrt{N_\lambda+1}}{|\alpha_\lambda|}
\right)
\int_{t_i}^{t_f}\dd t\,
\sum_\lambda\mathcal E_{\lambda,\rm cl}(t)
+\epsilon_{\rm cl,tail}.
\end{align*}
这就是正文\eqnrefs{eq:multimode_classical_limit_duhamel_bound} 的来源。它同时说明了经典场近似的三个独立要求：平均场幅在所研究时间窗内保持有限、量子涨落相对幅度随 $|\alpha_\lambda|^{-1}$ 变小，以及实际动力学不能大量泄漏出所控制的电子/Fock 工作窗。
\end{proof}


单模体积缩放给出最直观的例子。对保持无量纲模形固定的一族量子化体积，可把一量子场幅和电子矩阵元写成 $E_{1}^{\rm mode}(V)=C_E V^{-1/2}$、$G(V)=C_GV^{-1/2}$，其中 $C_E,C_G$ 与 $V$ 无关；保持经典场幅固定则要求 $N(V)=|\alpha(V)|^2=C_NV$。因此 $|\alpha|G=\sqrt{C_N}C_G=O(1)$，而固定涨落扇区中的量子修正相对于平均驱动为 $O(\sqrt{N_{\rm fl}+1}/|\alpha|)$，随 $V\to\infty$ 消失。这一极限明确区分了有限平均场与可忽略量子涨落两个物理尺度。

\subsection{Fock 块动力学}

纯态振幅层级适合描述确定的联合态，但一般实验还需要处理热态、混合光场以及偏迹后的电子约化态。因此从同一个线性哈密顿量出发，改用联合密度算符。把自由电子与玻色模式各自的自由演化吸收到相互作用绘景后，原来 Schrödinger 绘景中的电子耦合算符 $\hat{\mathcal G}_\lambda$ 变成
\[
\hat{\mathcal G}_{\lambda,I}(t)
\equiv
\ee^{+\ii\hat H_e t/\hbar}\hat{\mathcal G}_\lambda
\ee^{-\ii\hat H_e t/\hbar}\ee^{-\ii\omega_\lambda t}.
\]
下标 $I$ 表示相互作用绘景；Fourier 域对象继续按全文约定用波浪号区分。相互作用哈密顿量相应写成
\[
\hat H_{{\rm int},I}(t)
=\hbar\sum_\lambda
\left[\hat{\mathcal G}_{\lambda,I}(t)\hat a_\lambda
+\hat{\mathcal G}_{\lambda,I}^\dagger(t)\hat a_\lambda^\dagger\right].
\]
在不作弱耦合、单模或反冲截断时，联合传播子就是
\begin{equation}
{
\hat U_I(t,t_0)
=\mathcal T\exp\!\left[-\ii\int_{t_0}^{t}\dd t'\,
\frac{\hat H_{{\rm int},I}(t')}{\hbar}\right].}
\label{eq:general_multimode_time_ordered_propagator}
\end{equation}
Dyson 级数只是这个精确表达式的耦合阶展开；只有当不同时间生成元具有特殊闭合对易结构时，Magnus 展开才可能在有限阶闭合。

精确传播子~\eqnrefs{eq:general_multimode_time_ordered_propagator} 同样决定联合密度算符的演化。对任意混合初态，相互作用绘景的 Liouville--von Neumann 方程为
\begin{equation}
{
\ii\dot{\hat\rho}_I(t)
=\left[\frac{\hat H_{{\rm int},I}(t)}{\hbar},\hat\rho_I(t)\right].}
\label{eq:general_multimode_von_neumann_equation}
\end{equation}
为了显式保留环境布居与相干，先取多模 Fock 基
$|\vb{n}\rangle\equiv|n_1,n_2,\ldots\rangle$，其中 $n_\lambda=0,1,2,\ldots$ 是第 $\lambda$ 个玻色模式的占据数；再定义玻色 Fock 分块
\[
\hat\rho^{(e)}_{\vb{n}\vb{m}}(t)
\equiv
\langle\vb{n}|\hat\rho_I(t)|\vb{m}\rangle .
\]
每个分块仍是电子 Hilbert 空间上的算符；$\vb{n}=\vb{m}$ 描述环境占据数扇区，$\vb{n}\ne\vb{m}$ 保留环境相干性。把精确 Liouville--von Neumann 方程投影到这些分块，得到
\begin{equation}
{
\begin{aligned}
\ii\dot\hat\rho^{(e)}_{\vb{n}\vb{m}}
=\sum_\lambda\Big[&
\sqrt{n_\lambda+1}\,\hat{\mathcal G}_{\lambda,I}(t)
\hat\rho^{(e)}_{\vb{n}+\vb{e}_\lambda,\vb{m}}
+\sqrt{n_\lambda}\,\hat{\mathcal G}_{\lambda,I}^\dagger(t)
\hat\rho^{(e)}_{\vb{n}-\vb{e}_\lambda,\vb{m}}\\
&-\sqrt{m_\lambda}\,
\hat\rho^{(e)}_{\vb{n},\vb{m}-\vb{e}_\lambda}
\hat{\mathcal G}_{\lambda,I}(t)
-\sqrt{m_\lambda+1}\,
\hat\rho^{(e)}_{\vb{n},\vb{m}+\vb{e}_\lambda}
\hat{\mathcal G}_{\lambda,I}^\dagger(t)
\Big].
\end{aligned}}
\label{eq:universal_multiboson_density_block_equation}
\end{equation}
电子约化态只在最后对环境对角块求和，
\[
\hat\rho_e(t)=\sum_{\vb{n}}\hat\rho^{(e)}_{\vb{n}\vb{n}}(t).
\]
因此\eqnrefs{eq:universal_multiboson_density_block_equation} 比纯态振幅层级多保留了两类信息：初始或动态产生的环境混合，以及不同 Fock 扇区之间的相干。总联合动力学仍然幺正，所以全部对角分块的电子迹之和满足
\begin{equation}
{
\frac{\dd}{\dd t}
\sum_{\vb{n}}\operatorname{Tr}_e
\hat\rho^{(e)}_{\vb{n}\vb{n}}(t)=0.}
\label{eq:universal_multiboson_trace_conservation}
\end{equation}
偏迹后电子自身仍可以失去纯度，这正是电子--环境纠缠的结果，而不是总概率损失。

\begin{proposition}[Fock 层级的总迹守恒]
对完整玻色 Fock 格，若密度块及哈密顿量乘积的迹求和绝对收敛，\eqnrefs{eq:universal_multiboson_density_block_equation}给 $\frac{\dd}{\dd t}\sum_{\vb{n}}\Tr_e\rho^{(e)}_{\vb{n}\vb{n}}=0$。有限维自伴截断也保持迹，但完整格向截断格的泄漏需另估计。
\end{proposition}
\begin{proof}
在正文\eqnrefs{eq:universal_multiboson_density_block_equation} 中取 $\vb{m}=\vb{n}$，再对电子 Hilbert 空间取迹。对固定模式 $\lambda$，四个相邻块贡献为
\begin{align*}
\ii\frac{\dd}{\dd t}
\operatorname{Tr}_e\hat\rho^{(e)}_{\vb{n}\vb{n}}
=\sum_\lambda\operatorname{Tr}_e\Big[&
\sqrt{n_\lambda+1}\,
\hat{\mathcal G}_{\lambda,I}
\hat\rho^{(e)}_{\vb{n}+\vb{e}_\lambda,\vb{n}}
+\sqrt{n_\lambda}\,
\hat{\mathcal G}_{\lambda,I}^\dagger
\hat\rho^{(e)}_{\vb{n}-\vb{e}_\lambda,\vb{n}}\\
&-\sqrt{n_\lambda}\,
\hat\rho^{(e)}_{\vb{n},\vb{n}-\vb{e}_\lambda}
\hat{\mathcal G}_{\lambda,I}
-\sqrt{n_\lambda+1}\,
\hat\rho^{(e)}_{\vb{n},\vb{n}+\vb{e}_\lambda}
\hat{\mathcal G}_{\lambda,I}^\dagger
\Big].
\end{align*}
现在对全部 $\vb{n}$ 求和。第一项令 $\vb{r}=\vb{n}+\vb{e}_\lambda$，并用迹的循环性，得到
\begin{align*}
\sum_{\vb{n}}\sqrt{n_\lambda+1}\,
\operatorname{Tr}_e\!\left(
\hat{\mathcal G}_{\lambda,I}
\hat\rho^{(e)}_{\vb{n}+\vb{e}_\lambda,\vb{n}}
\right)
=\sum_{\vb{r}}\sqrt{r_\lambda}\,
\operatorname{Tr}_e\!\left(
\hat\rho^{(e)}_{\vb{r},\vb{r}-\vb{e}_\lambda}
\hat{\mathcal G}_{\lambda,I}
\right),
\end{align*}
它逐项抵消第三项。第二项令 $\vb{r}=\vb{n}-\vb{e}_\lambda$，同理有
\begin{align*}
\sum_{\vb{n}}\sqrt{n_\lambda}\,
\operatorname{Tr}_e\!\left(
\hat{\mathcal G}_{\lambda,I}^\dagger
\hat\rho^{(e)}_{\vb{n}-\vb{e}_\lambda,\vb{n}}
\right)
=\sum_{\vb{r}}\sqrt{r_\lambda+1}\,
\operatorname{Tr}_e\!\left(
\hat\rho^{(e)}_{\vb{r},\vb{r}+\vb{e}_\lambda}
\hat{\mathcal G}_{\lambda,I}^\dagger
\right),
\end{align*}
它逐项抵消第四项。因此每个 $\lambda$ 的总和都为零，最终得到
\[
\frac{\dd}{\dd t}\sum_{\vb{n}}\operatorname{Tr}_e
\hat\rho^{(e)}_{\vb{n}\vb{n}}=0,
\]
即正文\eqnrefs{eq:universal_multiboson_trace_conservation}。
\end{proof}


\begin{proposition}[线性玻色耦合的密度块邻接层级]
在上述共同域条件下，任意联合密度算符的 Fock 块满足\eqnrefs{eq:universal_multiboson_density_block_equation}。左、右指标分别按 ket 与 bra 阶梯矩阵元移动，Liouville 层级仅含每个 Fock 指标的最近邻块。
\end{proposition}
\begin{proof}
把相互作用绘景哈密顿量写成
\[
\hat H_I(t)
=\hbar\sum_\lambda
\left[
\hat{\mathcal G}_{\lambda,I}(t)\hat a_\lambda
+\hat{\mathcal G}_{\lambda,I}^\dagger(t)\hat a_\lambda^\dagger
\right],
\]
其中电子算符与 Fock 算符作用在不同张量因子上，彼此对易。定义
\[
\hat\rho^{(e)}_{\vb{n}\vb{m}}
\equiv\langle\vb{n}|\hat\rho_I|\vb{m}\rangle.
\]
左乘项的指标移动来自 ket 一侧。例如
\begin{align*}
\langle\vb{n}|\hat a_\lambda\hat\rho_I|\vb{m}\rangle
&=\sum_{\vb{r}}
\langle\vb{n}|\hat a_\lambda|\vb{r}\rangle
\hat\rho^{(e)}_{\vb{r}\vb{m}}\\
&=\sqrt{n_\lambda+1}\,
\hat\rho^{(e)}_{\vb{n}+\vb{e}_\lambda,\vb{m}},\\
\langle\vb{n}|\hat a_\lambda^\dagger\hat\rho_I|\vb{m}\rangle
&=\sqrt{n_\lambda}\,
\hat\rho^{(e)}_{\vb{n}-\vb{e}_\lambda,\vb{m}}.
\end{align*}
右乘项不能照抄左式，因为升降算符作用在 bra 指标上。插入完备关系后
\begin{align*}
\langle\vb{n}|\hat\rho_I\hat a_\lambda|\vb{m}\rangle
&=\sum_{\vb{r}}
\hat\rho^{(e)}_{\vb{n}\vb{r}}
\langle\vb{r}|\hat a_\lambda|\vb{m}\rangle\\
&=\sqrt{m_\lambda}\,
\hat\rho^{(e)}_{\vb{n},\vb{m}-\vb{e}_\lambda},\\
\langle\vb{n}|\hat\rho_I\hat a_\lambda^\dagger|\vb{m}\rangle
&=\sqrt{m_\lambda+1}\,
\hat\rho^{(e)}_{\vb{n},\vb{m}+\vb{e}_\lambda}.
\end{align*}
现在对
\[
\ii\hbar\dot\rho_I=\hat H_I\rho_I-\rho_I\hat H_I
\]
取 $\langle\vb{n}|\cdots|\vb{m}\rangle$。前两条左乘矩阵元带正号，后两条右乘矩阵元带负号；约去共同的 $\hbar$ 后得到正文\eqnrefs{eq:universal_multiboson_density_block_equation}。

这个推导还给出一个结构信息：刘维尔生成元在 Fock 指标 $(\vb{n},\vb{m})$ 上只连接每个模式的最近邻
\[
(\vb{n},\vb{m})
\longleftrightarrow
(\vb{n}\pm\vb{e}_\lambda,\vb{m}),
\quad
(\vb{n},\vb{m}\pm\vb{e}_\lambda),
\]
所以数值截断时边界误差来自被截掉的邻接块，而不是来自方程内部额外的非局域 Fock 跳跃。
\end{proof}


若联合态恰为纯态，写成
$\hat\rho^{(e)}_{\vb{n}\vb{m}}=|\psi^{(e)}_{\vb{n}}\rangle\langle\psi^{(e)}_{\vb{m}}|$，再投影到电子动量--自旋基，就恢复\eqnrefs{eq:universal_multiboson_amplitude_hierarchy}。因此振幅层级正是密度块方程限制到联合纯态流形后的严格子情形。

\subsection{有限时间跃迁与 FGR}

\eqnrefs{eq:universal_multiboson_amplitude_hierarchy} 给出完整有限时间动力学。弱耗尽时可截取 Dyson 级数的一阶振幅；作用时间进一步长于末态谱的相关时间时，一阶有限宽谱核收缩为 Fermi 黄金律。这里\emph{Fermi 黄金律}定义为一阶含时微扰概率在“弱跃迁 + 连续末态 + 足够长观测时间”共同极限下得到的恒定跃迁率：有限时间的 $\sinc^2$ 频谱选择核在分布意义下收缩为能量守恒 $\delta$ 函数。连续态归一化先由体积 $V$ 的周期盒调节，取单位归一电子本征态 $|i_e\rangle_V,|f_e\rangle_V$，能量分别为 $E_i,E_f$。若电子发射一个第 $\lambda$ 模量子，定义离散通道的一量子矩阵元与频率耦合
\[
M_{fi,\lambda}^{\rm em}
\equiv{}_V\!\langle f_e|\delta\hat h_\lambda^\dagger|i_e\rangle_V,
\qquad
\mathcal G_{fi,\lambda}^{\rm em}=M_{fi,\lambda}^{\rm em}/\hbar,
\]
其中 $[\mathcal G^{\rm em}]={\rm s^{-1}}$。初始模式占据数为 $n_\lambda$ 时，一阶有限时间发射概率为
\begin{equation}
{
P_{i\to f}^{\rm em}(T)
=(n_\lambda+1)
|\mathcal G_{fi,\lambda}^{\rm em}|^2T^2
\sinc^2\!\left(\frac{\Delta_{fi}^{\rm em}T}{2}\right),
\qquad
\Delta_{fi}^{\rm em}
=\frac{E_f-E_i}{\hbar}+\omega_\lambda .}
\label{eq:universal_boson_finite_time_emission}
\end{equation}
$\Delta_{fi}^{\rm em}$ 的 SI 单位是 ${\rm s^{-1}}$，它比较电子损失的角频率与模式角频率。这里 $\sinc x\equiv\sin x/x$；有限作用时间 $T$ 因而把严格能量守恒展宽成宽度为 $O(T^{-1})$ 的频率选择窗。$n_\lambda+1$ 同时包含真空自发发射以及已有玻色占据产生的受激增强。对应吸收通道使用 $\delta\hat h_\lambda$，定义
\[
M_{fi,\lambda}^{\rm abs}
\equiv{}_V\!\langle f_e|\delta\hat h_\lambda|i_e\rangle_V,
\qquad
\mathcal G_{fi,\lambda}^{\rm abs}=M_{fi,\lambda}^{\rm abs}/\hbar.
\]
其概率为
\[
P_{i\to f}^{\rm abs}(T)
=n_\lambda|\mathcal G_{fi,\lambda}^{\rm abs}|^2T^2
\sinc^2\!\left(\frac{\Delta_{fi}^{\rm abs}T}{2}\right),
\qquad
\Delta_{fi}^{\rm abs}=\frac{E_f-E_i}{\hbar}-\omega_\lambda,
\]
所以真空不存在对应的吸收支路。

有限时间的 $\sinc^2$ 线宽为 $O(T^{-1})$。只有当末态形成连续谱，而且矩阵元和态密度在这一共振窗内变化缓慢时，才可以使用分布极限
\begin{equation}
{
T\sinc^2\!\left(\frac{\delta E\,T}{2\hbar}\right)
\xrightarrow[T\to\infty]{\mathcal D'}
2\pi\hbar\,\delta(\delta E).}
\label{eq:finite_time_sinc_distribution_limit}
\end{equation}
这里 $\delta E$ 是能量失配。窄离散模、强耦合、高占据数、显著初态耗尽或阈值附近快速变化的态密度都会破坏这一极限，此时应继续使用有限时间振幅层级。

\begin{lemma}[有限时间 sinc 平方核的分布极限]
令 $D_T(\Omega)=T\operatorname{sinc}^2(\Omega T/2)$，$T>0$。则 $D_T\to2\pi\delta$ 于分布空间。对一般色散失谐 $\Omega(\vb{\kappa})$，只有在零能壳为正则水平集且耦合加权后的推前谱密度足够光滑可积时，才可把该极限代入谱积分，得到本章黄金律；阈值奇点须保留有限时间核。
\end{lemma}
\begin{proof}
$T\to\infty$ 始终按分布意义理解。 长时间分布极限只在作用时间远大于相关时间、且能谱在有限时间核的特征能宽 $\Delta E_T\equiv 2\hbar/T$ 内可视为平滑时使用。 先定义
$D_T(\Omega)=T\sinc^2(\Omega T/2)$。 对任意光滑紧支撑测试函数 $F$，令 $x=\Omega T/2$，则

\[
\int\dd\Omega\,F(\Omega)D_T(\Omega)
=2\int\dd x\,F(2x/T)\sinc^2x.
\]
当 $T\to\infty$ 时，$F(2x/T)\to F(0)$，且
$\int_{-\infty}^{\infty}\sinc^2x\,\dd x=\pi$，所以
\[
D_T(\Omega)\xrightarrow{\mathcal D'}2\pi\delta(\Omega).
\]
再用 $\Omega=\delta E/\hbar$ 与
$\delta(\delta E/\hbar)=\hbar\delta(\delta E)$，得到相应的能量形式。分布极限的支配收敛由 $|F(2x/T)\operatorname{sinc}^2x|\le\|F\|_\infty\operatorname{sinc}^2x$ 给出。核归一化也可由矩形脉冲的 Fourier Parseval 直接得到 $\int\operatorname{sinc}^2x\dd x=\pi$。
\end{proof}


\begin{proposition}[有限时间谱核的归一化与黄金律时间窗]
把所有末态的失谐写为 $\delta$（单位 $\si{\per\second}$），并把耦合平方与末态测度合并成非负密度
$J(\delta)$，使一阶总跃迁概率为
\begin{equation}
P^{(2)}(T)=\int_{\mathbb R}J(\delta)T^2
\sinc^2\!\left(\frac{\delta T}{2}\right)\,\dd\delta.
\label{eq:feqo-finite-time-spectral-density-probability}
\end{equation}
若 $J$ 有界且在 $\delta=0$ 连续，则
$P^{(2)}(T)/T\to2\pi J(0)$。
这个极限只是二阶表达式的谱极限；用于真实动力学还要求该时间窗内初态未显著耗尽。
\end{proposition}
\begin{proof}
令 $K_T(\delta)=T\sinc^2(\delta T/2)/(2\pi)$。
它非负，且由 $\int_{\mathbb R}(\sin x/x)^2\dd x=\pi$ 得
$\int K_T\dd\delta=1$。对任意 $\Omega>0$，
\[
\int_{|\delta|>\Omega}K_T(\delta)\dd\delta
\le\frac{4}{\pi T\Omega},
\]
因为 $\sin^2\le1$。若 $\omega_J(\Omega)
=\sup_{|\delta|\le\Omega}|J(\delta)-J(0)|$，则
\[
\left|\frac{P^{(2)}(T)}T-2\pi J(0)\right|
\le2\pi\omega_J(\Omega)+\frac{16\|J\|_\infty}{T\Omega}.
\]
先固定足够小的 $\Omega$，再令 $T\Omega$ 足够大，即得到收敛。
这里没有把离散末态变成连续谱：若 $J$ 含孤立 delta 质量，以上连续性条件不成立。
\end{proof}

\begin{exbox}[从短时二次增长到恒定速率：一个可积谱核]
设电子投影后的末态失谐谱在共振附近可用
\[
J(\delta)=\frac{g_*^2a}{\pi(\delta^2+a^2)},\qquad a>0,
\]
表示，其中 $g_*,a$ 都是角频率尺度。失谐变量可正可负；该洛伦兹密度是工作窗模型，
不是把环境的绝对频率延伸到负值。它的相关核由 Fourier 积分给出
$C(\tau)=\int J(\delta)\ee^{\ii\delta\tau}\dd\delta
=g_*^2\ee^{-a|\tau|}$。
将有限时间振幅平方写成两个时间积分，得到
\begin{align*}
P^{(2)}(T)&=\int_0^T\dd t\int_0^T\dd s\,C(t-s)\\
&=2g_*^2\int_0^T(T-\tau)\ee^{-a\tau}\dd\tau\\
&=\frac{2g_*^2}{a^2}(aT-1+\ee^{-aT}).
\end{align*}
于是
\begin{equation}
P^{(2)}(T)=
\begin{cases}
g_*^2T^2[1-aT/3+O((aT)^2)],&aT\ll1,\\
\Gamma_*T-\Gamma_*/a+O(\ee^{-aT}),&aT\gg1,
\end{cases}
\qquad \Gamma_*=\frac{2g_*^2}{a}.
\label{eq:feqo-lorentzian-spectrum-golden-window}
\end{equation}
黄金律需要 $a^{-1}\ll T\ll\Gamma_*^{-1}$，所以存在宽时间窗还要求 $g_*\ll a$。
例如取 $g_*/a=0.05$、$aT=20$，得到 $P^{(2)}\simeq0.095$，
黄金律给 $\Gamma_*T=0.100$；两者相差 $0.005$，正是有限建立时间造成的修正。
若窄共振使 $g_*\gtrsim a$，上述两侧时间尺度不再分离，便应保留环境模的联合动力学。
\end{exbox}


\begin{proposition}[连续频率共振根的雅可比因子]
设连续模按角频率归一化，$[a(\omega),a^\dagger(\omega')]=\delta(\omega-\omega')$。
将选定电子吸收通道的矩阵元并入 $g(\omega)$，使吸收项写为
$\hbar\int g(\omega)B^\dagger a(\omega)\dd\omega$，
其中 $B^\dagger$ 是该通道的无量纲电子跃迁算符，$[g]={\rm s^{-1/2}}$。
取平稳且频率对角的光输入，满足
$\langle a^\dagger(\omega)a(\omega')\rangle=n(\omega)\delta(\omega-\omega')$，
$n$ 无量纲。令 $\Delta(\omega)$ 为实际吸收通道的角频率失谐，
则有限时间一阶吸收数谱为
\begin{equation}
 \frac{\dd P_+(T)}{\dd\omega}=|g(\omega)|^2 n(\omega)
 T^2\sinc^2[\Delta(\omega)T/2].
 \label{eq:feqo-continuum-mode-amplitude}
\end{equation}
若权重 $|g|^2n$ 具有紧支撑且连续，$\Delta$ 在支持区内连续可微，
共振根 $\omega_j$ 有限、位于频率域内部且满足 $\Delta'(\omega_j)\ne0$，则
\begin{equation}
 \lim_{T\to\infty}\frac{P_+(T)}T=
 \Gamma_+=2\pi\sum_j\frac{|g(\omega_j)|^2n(\omega_j)}{|\Delta'(\omega_j)|}.
 \label{eq:feqo-continuum-absorption}
\end{equation}
这里仍是弱跃迁表达式的谱极限，真实动力学还需满足未显著耗尽的时间窗条件。
\end{proposition}
\begin{proof}
一阶吸收振幅正比于 $g(\omega)\int_0^T\ee^{\ii\Delta(\omega)t}\dd t$。
对光输入求期望时，频率对角相关使两个频率积分收缩为一个，
留下占据因子 $n(\omega)$，得到所给数谱。
左边单位为秒，和右边的 $|g|^2T^2$ 一致。

每个简单根附近由反函数定理作变量替换 $u=\Delta(\omega)$，
并记该支的反函数为 $\omega_j(u)$。局部概率除以时间后成为
\[
\int\frac{|g(\omega_j(u))|^2n(\omega_j(u))}
{|\Delta'(\omega_j(u))|}\,
T\sinc^2(uT/2)\dd u.
\]
推前权重在零点连续，所以有限时间核的近似恒等式给出
$2\pi|g(\omega_j)|^2n(\omega_j)/|\Delta'(\omega_j)|$。
剩余紧支撑区域与各根分离，失谐绝对值有正下界；
利用 $T\sinc^2(\Delta T/2)\le4/(T\Delta^2)$，
这一部分随 $T^{-1}$ 消失。将各支贡献相加，便得所给速率。
\end{proof}

无共振根时，上述速率为零。
多重根或带边处的雅可比因子发散表示简单根黄金律失效，
应保留有限时间核与真实谱结构，不能把发散速率当作实验概率。

\begin{exbox}[真空单光子通道是否存在]
考虑无外部结构、无其他粒子的自由电子和一个真实真空光子。
检验电子单独吸收或发射该光子能否同时满足能量和动量守恒，
并说明连续模式求和中的体积因子抵消是否足以保证非零跃迁率。

\textbf{解。} 真空光子四波矢为 $K^\mu=(\omega/c,\vb{k})$，
$K^2=0$、$\omega>0$；自由电子满足 $P^2=m_{\rm e}^2c^2$。
若电子吸收一个真实光子，末态四动量为 $P+\hbar K$，
质量壳要求
\[
(P+\hbar K)^2=P^2\quad\Longrightarrow\quad 2\hbar P\cdot K=0.
\]
然而电子静止系中 $P\cdot K=m_{\rm e}\omega>0$，且此内积为 Lorentz 不变量，矛盾。
发射的末态 $P-\hbar K$ 也要求同一不可能条件。
因此这一渐近通道没有同时满足全部守恒条件的共振根。
盒归一化中模式态密度的体积因子与矩阵元平方确会抵消，
但这种形式操作并不会创造一个原本不存在的质量壳交点。
近场、介质、周期结构或有物理开关的外源能提供额外动量与能量，
这才是前面有限窗口及结构谱能够产生真实信号的条件。
\end{exbox}


对连续正常模还必须把归一化约定与模式态密度一起处理。令 $\dd\mu_\nu(\vb{\kappa})$ 是分支 $\nu$ 的模式测度，并定义振幅密度 $\mathfrak g_{\nu\vb{G}}(\vb{\kappa})$，使得
\[
|\mathfrak g_{\nu\vb{G}}(\vb{\kappa})|^2\dd\mu_\nu(\vb{\kappa})
\equiv
\lim_{V\to\infty}
\sum_{\vb{\kappa}_{\vb{n}}\in\dd\mu_\nu(\vb{\kappa})}
|\mathcal G_{fi}^{(\nu,V)}(\vb{\kappa}_{\vb{n}},\vb{G})|^2 .
\]
$\mathfrak g$ 单独的量纲依赖于测度约定，而组合
$|\mathfrak g|^2\dd\mu$ 总具有 ${\rm s^{-2}}$ 量纲。对于初始电子动量 $\vb{p}_i$，向模式 $(\nu,\vb{\kappa})$ 发射一个量子并允许倒格矢 $\vb{G}$ 时，末态动量为
$\vb{p}_f=\vb{p}_i-\hbar(\vb{\kappa}+\vb{G})$，故精确能量失配是
\[
\delta E_{\nu\vb{\kappa},\vb{G}}
=E[\vb{p}_i-\hbar(\vb{\kappa}+\vb{G})]-E(\vb{p}_i)
+\hbar\omega_{\nu\vb{\kappa}}.
\]
把\eqnrefs{eq:finite_time_sinc_distribution_limit} 用于有限时间概率测度，得到连续正常模的 Fermi 黄金规则速率测度
\begin{equation}
{
\dd\Gamma_{\nu\vb{G}}^{\rm em}
=2\pi(n_{\nu\vb{\kappa}}+1)
|\mathfrak g_{\nu\vb{G}}(\vb{\kappa})|^2
\delta\!\left[
\frac{E(\vb{p}_i)-E[\vb{p}_i-\hbar(\vb{\kappa}+\vb{G})]}{\hbar}
-\omega_{\nu\vb{\kappa}}
\right]\dd\mu_\nu(\vb{\kappa}).}
\label{eq:general_free_electron_golden_rule}
\end{equation}
\begin{corollary}[连续通道的黄金律速率测度]
若失谐的零水平集正则，且耦合加权推前密度满足前述连续性与可积条件，
则一阶概率除以时间的长时间极限为\eqnrefs{eq:general_free_electron_golden_rule}。
该结论作用于弱跃迁的谱表达式；有限耗尽条件仍须独立检验。
\end{corollary}
\begin{proof}
连续模式的有限时间概率测度是
\[
\dd P^{\rm em}
=(n_{\nu\vb{\kappa}}+1)|\mathfrak g_{\nu\vb{G}}|^2T^2
\sinc^2\!\left(\frac{\delta E T}{2\hbar}\right)
\dd\mu_\nu(\vb{\kappa}).
\]
把它除以作用时间 $T$，并把能量失配改写成角频率失配
\[
\Omega_{\nu\vb{G}}(\vb{\kappa})
\equiv\frac{\delta E_{\nu\vb{\kappa},\vb{G}}}{\hbar},
\]
可得
\begin{align*}
\frac{\dd P^{\rm em}}{T}
&=(n_{\nu\vb{\kappa}}+1)|\mathfrak g_{\nu\vb{G}}|^2
\left[T\sinc^2\!\left(\frac{\Omega_{\nu\vb{G}}T}{2}\right)\right]
\dd\mu_\nu(\vb{\kappa})\\
&\xrightarrow[T\to\infty]{}
2\pi(n_{\nu\vb{\kappa}}+1)|\mathfrak g_{\nu\vb{G}}|^2
\delta\!\left[\Omega_{\nu\vb{G}}(\vb{\kappa})\right]
\dd\mu_\nu(\vb{\kappa}).
\end{align*}
发射末态满足
\[
\vb{p}_f=\vb{p}_i-\hbar(\vb{\kappa}+\vb{G}),
\]
所以
\begin{align*}
\Omega_{\nu\vb{G}}(\vb{\kappa})
&=\frac{E(\vb{p}_f)-E(\vb{p}_i)}{\hbar}
+\omega_{\nu\vb{\kappa}}\\
&=-\left\{
\frac{E(\vb{p}_i)-E[\vb{p}_i-\hbar(\vb{\kappa}+\vb{G})]}{\hbar}
-\omega_{\nu\vb{\kappa}}
\right\}.
\end{align*}
利用 $\delta(-x)=\delta(x)$，得到\eqnrefs{eq:general_free_electron_golden_rule}。量纲在这里同时闭合：
$|\mathfrak g|^2\dd\mu$ 的量纲为 ${\rm s^{-2}}$，角频率 $\delta$ 分布的量纲为 ${\rm s}$，故 $\dd\Gamma$ 的量纲为 ${\rm s^{-1}}$。若改用能量形式，则
\[
\delta(\Omega)=\delta(\delta E/\hbar)
=\hbar\,\delta(\delta E),
\]
这就是能量表象黄金律中相应 $\hbar$ 因子的来源。

\end{proof}

只有对实际允许的 $\vb{\kappa}$、$\nu$、$\vb{G}$ 和末态自旋积分/求和以后，才得到单位为 ${\rm s^{-1}}$ 的总率。这里始终使用精确相对论色散
$E(\vb{p})=\sqrt{m_{\mathrm e}^2c^4+c^2\vb{p}^2}$；低反冲同步关系必须在此之后作为受控展开取得。

令 $\vb{v}_i=\nabla_{\vb{p}}E|_{\vb{p}_i}$，$H_E$ 是 $E(\vb{p})$ 对动量的海森矩阵。若实际支持区满足二阶及更高反冲频移均小于所需分辨尺度，则
\[
\frac{E(\vb{p}_i)-E[\vb{p}_i-\hbar(\vb{\kappa}+\vb{G})]}{\hbar}
=\vb{v}_i\cdot(\vb{\kappa}+\vb{G})
-\frac{\hbar}{2}(\vb{\kappa}+\vb{G})^{\mathsf T}
H_E(\vb{p}_i)(\vb{\kappa}+\vb{G})+\cdots,
\]
最低阶才化为
\begin{equation}
{
\omega_{\nu\vb{\kappa}}\simeq
\vb{v}_i\cdot(\vb{\kappa}+\vb{G}).}
\label{eq:generalized_cherenkov_condition}
\end{equation}
$\vb{G}=0$ 且模式是均匀透明介质的传播电磁模时，这是 Cherenkov 型同步；离散平移对称系统允许 $\vb{G}\neq0$ 时，同一关系包含 Smith--Purcell 型倒格矢补偿。二者都是精确能量守恒在小反冲展开中的最低阶同步条件。\cite{RiveraKaminer2020Quasiparticles,GarciaDeAbajo2010RMP}

有限时间跃迁率只保留了一阶弱跃迁的信息。若实验真正记录的是电子末态能谱，更一般的对象是电子能量改变的完整概率分布。先取电子自由哈密顿量的谱投影 $\Pi_E$；为避免连续谱的分布归一化问题，可以先在有限盒或有限能量分辨率下离散化，最后再取谱积分极限。设初态电子已经在能量测量意义下去相干，
\[
\rho_\ee^{\rm diag}=\sum_n p_n|n\rangle\langle n|,
\qquad
H_e|n\rangle=E_n|n\rangle,
\]
环境初态为任意 $\rho_X$，联合有限时间传播子为 $U(T)$。定义电子能量改变
\[
\Delta E\equiv E_f-E_i,
\]
于是两次测量（TPM）能量交换分布为
\begin{equation}
{
P_T(\Delta E)
=\sum_{m,n}p_n\,
\operatorname{Tr}_X\!\left[
\langle m|U(T)|n\rangle\rho_X
\langle n|U^\dagger(T)|m\rangle
\right]
\delta[\Delta E-(E_m-E_n)].}
\label{eq:electron_energy_exchange_tpm_distribution}
\end{equation}
它已经包含任意阶多量子过程、反冲和环境末态求和；FGR 只是它在弱耦合、长时间并只保留最低非零阶后的局部近似。这里约定电子能量增加为 $\Delta E>0$；EELS 常用的“损失能量”则是 $E_{\rm loss}=-\Delta E>0$。

比直接储存整个分布更方便的是特征函数。引入有时间量纲的计数场 $\lambda$，
\begin{equation}
{
\chi_E(\lambda)
\equiv\int\dd(\Delta E)\,
P_T(\Delta E)\ee^{+\ii\lambda\Delta E/\hbar}
=\operatorname{Tr}_{eX}\!\left[
\ee^{+\ii\lambda H_e/\hbar}
U(T)
\ee^{-\ii\lambda H_e/\hbar}
(\rho_\ee^{\rm diag}\otimes\rho_X)
U^\dagger(T)
\right].}
\label{eq:electron_energy_exchange_characteristic_function}
\end{equation}
归一性给 $\chi_E(0)=1$。电子能量改变的累积量由
\[
\mathcal C_n^{(E)}
=\left(\frac{\hbar}{\ii}\right)^n
\left.\frac{\partial^n}{\partial\lambda^n}
\ln\chi_E(\lambda)\right|_{\lambda=0}
\]
统一给出；$\mathcal C_1^{(E)}$ 是平均电子能量改变，$\mathcal C_2^{(E)}$ 是能量扩散方差，更多阶则区分非高斯能损统计。

\begin{proposition}[两次能量测量的特征函数]
设初态因子化，先对电子作完整谱能量测量，再经联合幺正演化和末能量测量。初态经首次测量去相干为 $\rho_e^{\rm diag}$，能量差分布归一且非负，其特征函数为\eqnrefs{eq:electron_energy_exchange_characteristic_function}。能量有简并时使用谱投影块去相干；连续测量按谱测度解释。
\end{proposition}
\begin{proof}
初始能量测量得到结果 $E_n$ 的概率为 $p_n$，测量后联合态为
\[
|n\rangle\langle n|\otimes\rho_X.
\]
经过 $U(T)$ 后，再测得电子能量 $E_m$ 的条件概率是
\[
p(m|n)
=\operatorname{Tr}_{eX}\!\left[
(|m\rangle\langle m|\otimes\mathbb I_X)
U(|n\rangle\langle n|\otimes\rho_X)U^\dagger
\right].
\]
在电子空间取矩阵元，得到
\[
p(m|n)
=\operatorname{Tr}_X\!\left[
\langle m|U|n\rangle\rho_X
\langle n|U^\dagger|m\rangle
\right].
\]
因此联合概率 $p_n p(m|n)$ 按
$\Delta E=E_m-E_n$ 分组以后，正是正文\eqnrefs{eq:electron_energy_exchange_tpm_distribution}。

对该分布作 Fourier 变换：
\begin{align*}
\chi_E(\lambda)
&=\sum_{mn}p_n p(m|n)
\ee^{+\ii\lambda(E_m-E_n)/\hbar}\\
&=\sum_{mn}p_n
\operatorname{Tr}_X\!\left[
\ee^{+\ii\lambda E_m/\hbar}
\langle m|U|n\rangle
\ee^{-\ii\lambda E_n/\hbar}
\rho_X
\langle n|U^\dagger|m\rangle
\right].
\end{align*}
利用
\[
\ee^{+\ii\lambda H_e/\hbar}|m\rangle
=\ee^{+\ii\lambda E_m/\hbar}|m\rangle,
\qquad
\ee^{-\ii\lambda H_e/\hbar}|n\rangle
=\ee^{-\ii\lambda E_n/\hbar}|n\rangle,
\]
并在 $m,n$ 上插入完备关系，和式重新组合成
\[
\chi_E(\lambda)
=\operatorname{Tr}_{eX}\!\left[
\ee^{+\ii\lambda H_e/\hbar}
U
\ee^{-\ii\lambda H_e/\hbar}
(\rho_\ee^{\rm diag}\otimes\rho_X)
U^\dagger
\right],
\]
得到正文\eqnrefs{eq:electron_energy_exchange_characteristic_function}。

令 $\lambda=0$，幺正性给
\[
\chi_E(0)
=\operatorname{Tr}[U(\rho_\ee^{\rm diag}\otimes\rho_X)U^\dagger]=1.
\]
对 $\lambda$ 求导则有
\[
\left.\partial_\lambda\ln\chi_E\right|_0
=\frac{\ii}{\hbar}\langle\Delta E\rangle,
\]
继续求高阶导数便得到累积量生成公式。若能量改变只取离散值
$\ell\hbar\omega$，把 $\vartheta=\omega\lambda$ 代入即可得到正文\eqnrefs{eq:sideband_full_counting_characteristic}；所以边带特征函数并不是 PINEM 特有技巧，而是一般能量交换全计数统计的离散化。
\end{proof}


单频、低反冲平移格只是这套连续能量统计的离散特例。若所有可分辨末态满足
$\Delta E=\ell\hbar\omega$，定义无量纲计数角 $\vartheta=\omega\lambda$，便有
\begin{equation}
{
\chi_{\rm sb}(\vartheta)
=\sum_{\ell\in\ZZ}P_\ell\ee^{+\ii\ell\vartheta},
\qquad
P_\ell
=\int_{-\pi}^{\pi}\frac{\dd\vartheta}{2\pi}
\chi_{\rm sb}(\vartheta)\ee^{-\ii\ell\vartheta}.}
\label{eq:sideband_full_counting_characteristic}
\end{equation}
因此 PINEM 的 Bessel 边带、QPINEM 的 Fock 条件分布以及热光场产生的修正 Bessel 分布，都只是同一个生成函数在不同环境态和不同有效哈密顿量下的具体形式。

TPM 定义还提醒一个重要限制：如果入射电子本身含不同自由能量之间的相干，第一次理想能量测量会把这些相干去掉，所以\eqnrefs{eq:electron_energy_exchange_tpm_distribution} 描述的是“真正执行两次能量测量”的统计。若目标是利用入射谱相干做相位敏感 PINEM 或层析，就不能把第一次能量投影当成无害步骤；此时应直接传播完整密度矩阵，或使用允许准概率的相干全计数统计。能量布居与能量相干因此必须作为两类不同的测量信息处理。

TPM 特征函数还可以与开放系统章节的非高斯影响泛函直接连接。为此不再先对末态能量求和，而把计数场放进两条 Keldysh 支路上的电子传播子。定义对称倾斜（倾斜）传播子
\begin{equation}
{
\hat U_{\lambda/2}(T)
\equiv
\ee^{+\ii\lambda\hat H_e/(2\hbar)}
\hat U(T)
\ee^{-\ii\lambda\hat H_e/(2\hbar)},}
\label{eq:energy_counting_tilted_propagator}
\end{equation}
则当初态与 $\hat H_e$ 对易时，
\begin{equation}
{
\chi_E(\lambda)
=
\Tr_{eX}\!\left[
\hat U_{\lambda/2}(T)
(\rho_\ee^{\rm diag}\!\otimes\rho_X)
\hat U_{-\lambda/2}^{\dagger}(T)
\right].}
\label{eq:energy_counting_keldysh_trace}
\end{equation}
这与\eqnrefs{eq:electron_energy_exchange_characteristic_function} 完全等价，但更适合做微扰和环境累积量展开：正、反时间支路分别带 $+\lambda/2$ 与 $-\lambda/2$，电子每一次改变自由能量的顶点都因此获得计数相位。

若相互作用写成一般双线性形式
\[
\hat V_I(t)=\sum_a \hat S_a(t)\otimes\hat B_a(t),
\]
把带计数场的电子算符记为
\[
\hat S_{a,\lambda}(t)
=
\ee^{+\ii\lambda\hat H_e/(2\hbar)}
\hat S_a(t)
\ee^{-\ii\lambda\hat H_e/(2\hbar)},
\]
则特征函数可写成闭合时间路径生成泛函
\begin{equation}
{
\chi_E(\lambda)
=
\left\langle
\mathcal T_{\mathcal C}
\exp\!\left[
-\frac{\ii}{\hbar}
\int_{\mathcal C}\dd z\,
\sum_a
\hat S_{a,\lambda_{\mathcal C}}(z)\hat B_a(z)
\right]
\right\rangle .}
\label{eq:energy_fcs_contour_generating_functional}
\end{equation}
其中前向、后向支路的 $\lambda_{\mathcal C}$ 分别取 $+\lambda/2$、$-\lambda/2$。对环境做连通累积量展开后，
\begin{equation}
{
\ln\chi_E(\lambda)
=
\sum_{n=1}^{\infty}
\frac{(-\ii/\hbar)^n}{n!}
\int_{\mathcal C}\dd z_1\cdots\dd z_n\,
\big\langle\!\big\langle
\hat{\mathcal V}_{\lambda}(z_1)\cdots
\hat{\mathcal V}_{\lambda}(z_n)
\big\rangle\!\big\rangle_c,}
\label{eq:energy_fcs_connected_cumulant_expansion}
\end{equation}
这里 $\hat{\mathcal V}_{\lambda}\equiv\sum_a\hat S_{a,\lambda_{\mathcal C}}\hat B_a$。这条式子把“环境是否高斯”与“电子能损分布是否高斯”严格区分开：高斯环境只意味着环境连通相关函数在二阶终止，并不意味着电子的 $P_T(\Delta E)$ 必为高斯；即使环境是单模相干态或真空，多次离散量子交换仍可产生强烈非高斯的 Bessel、Poisson 或 Fock 条件分布。反过来，环境的三阶、四阶连通相关函数会直接进入 $\ln\chi_E$ 的更高阶项，使能量交换偏度、峰度以及不同交换事件之间的关联获得独立信息。

\begin{proposition}[计数倾斜与轮廓连通生成函数]
在两次能量测量的去相干初态下，$[\rho_e^{\rm diag},H_e]=0$，故倾斜传播子的迹满足\eqnrefs{eq:energy_counting_keldysh_trace}，等价于闭时轮廓生成泛函。其对数在源的解析非零邻域具有\eqnrefs{eq:energy_fcs_connected_cumulant_expansion}的连通展开；离开收敛域时只可作形式或截断展开。高斯环境不使非交换电子能量交换分布自动高斯。
\end{proposition}
\begin{proof}
先证明倾斜传播子的迹公式。由
\[
[\rho_\ee^{\rm diag},H_e]=0
\]
可把\eqnrefs{eq:electron_energy_exchange_characteristic_function} 中的
$\ee^{-\ii\lambda H_e/\hbar}$ 平分到 $\rho_\ee^{\rm diag}$ 两侧：
\begin{align*}
\chi_E(\lambda)
&=\Tr\!\left[
\ee^{+\ii\lambda H_e/\hbar}
U
\ee^{-\ii\lambda H_e/\hbar}
\rho\,U^\dagger
\right]\\
&=\Tr\!\left[
\ee^{+\ii\lambda H_e/(2\hbar)}
U
\ee^{-\ii\lambda H_e/(2\hbar)}
\rho\,
\ee^{-\ii\lambda H_e/(2\hbar)}
U^\dagger
\ee^{+\ii\lambda H_e/(2\hbar)}
\right]\\
&=\Tr\!\left[
U_{\lambda/2}\rho\,U_{-\lambda/2}^\dagger
\right],
\end{align*}
其中 $\rho=\rho_\ee^{\rm diag}\otimes\rho_X$，得到正文\eqnrefs{eq:energy_counting_keldysh_trace}。

再转到相互作用绘景。对前向支路，
\[
U_{\lambda/2,I}
=
\mathcal T
\exp\!\left[
-\frac{\ii}{\hbar}
\int_{t_i}^{t_f}\dd t\,
V_{I,+\lambda/2}(t)
\right],
\]
而后向支路
\[
U_{-\lambda/2,I}^{\dagger}
=
\widetilde{\mathcal T}
\exp\!\left[
+\frac{\ii}{\hbar}
\int_{t_i}^{t_f}\dd t\,
V_{I,-\lambda/2}(t)
\right].
\]
把两条指数放到同一闭合时间轮廓 $\mathcal C$ 上，后向支路的积分方向自动提供相反号，于是
\[
\chi_E(\lambda)
=
\left\langle
\mathcal T_{\mathcal C}
\exp\!\left[
-\frac{\ii}{\hbar}
\int_{\mathcal C}\dd z\,
V_{I,\lambda_{\mathcal C}}(z)
\right]
\right\rangle ,
\]
得到\eqnrefs{eq:energy_fcs_contour_generating_functional}。

最后使用任意轮廓算符的累积量生成恒等式
\[
\ln\left\langle
\mathcal T_{\mathcal C}\ee^{\hat X}
\right\rangle
=
\sum_{n=1}^{\infty}
\frac1{n!}
\left\langle\!\left\langle
\mathcal T_{\mathcal C}\hat X^n
\right\rangle\!\right\rangle_c,
\]
取
\[
\hat X
=
-\frac{\ii}{\hbar}
\int_{\mathcal C}\dd z\,
V_{I,\lambda_{\mathcal C}}(z),
\]
逐次展开即得正文\eqnrefs{eq:energy_fcs_connected_cumulant_expansion}。若环境算符 $B_a$ 处于高斯态且耦合对 $B_a$ 线性，环境迹中的 $n>2$ 连通相关函数消失；但电子算符仍保留时间排序和非交换性，重求和后的 $\chi_E(\lambda)$ 一般仍包含任意高阶 $\lambda$ 依赖。因此“高斯环境”与“高斯能量损失分布”在数学上是两个不同命题。
\end{proof}


因此若实验能够可靠测得电子能量改变的前三、四阶累积量，仅用 $\operatorname{Im}G^R$ 或二点噪声谱拟合并不总是充分；只有在环境高斯、耦合线性且电子动力学约化仍覆盖所有可分辨交换通道时，二点 Green/FDT 数据才封闭环境统计。非高斯缺陷、强非谐材料模或条件制备光场需要回到\eqnrefs{eq:energy_fcs_connected_cumulant_expansion} 与开放系统章节的高阶累积量层级。

为了与实验谱线统计直接连接，记电子能量改变的中心矩为 $\mu_n$。累积量与通常的偏度、峰度满足
\begin{equation}
{
\mu_3=\mathcal C_3^{(E)},
\qquad
\mu_4=\mathcal C_4^{(E)}+3[\mathcal C_2^{(E)}]^2,
\qquad
\operatorname{Skew}_E=\frac{\mathcal C_3^{(E)}}{[\mathcal C_2^{(E)}]^{3/2}},
\qquad
\operatorname{Kurt}_{E}^{\rm exc}=\frac{\mathcal C_4^{(E)}}{[\mathcal C_2^{(E)}]^2}.}
\label{eq:energy_fcs_skewness_kurtosis}
\end{equation}
其中 $\operatorname{Kurt}_{E}^{\rm exc}$ 是超额峰度。若探测器只叠加一个与真实电子能损独立的加性读出噪声 $N$，卷积在特征函数空间变成乘法，故累积量生成函数相加：
\begin{equation}
{
K_{\rm meas}(\lambda)
=K_E(\lambda)+K_N(\lambda),
\qquad
\mathcal C_n^{\rm meas}
=\mathcal C_n^{(E)}+\mathcal C_n^{(N)}.}
\label{eq:energy_fcs_detector_cumulant_addition}
\end{equation}
因此零均值高斯谱仪展宽只增加二阶累积量，不会伪造三、四阶连通累积量；但由于总方差变大，标准化的 $\operatorname{Skew}_E,\operatorname{Kurt}_{E}^{\rm exc}$ 会被稀释。这个区分使高阶 FCS 可以与普通仪器卷积模型直接兼容，而不需要先对测得谱线做不稳定的逐点反卷积。

单频离散格中还有简单尺度关系：若 $\Delta E=\ell\hbar\omega$，则
\begin{equation}
{
\mathcal C_n^{(E)}=(\hbar\omega)^n\mathcal C_n^{(\ell)}.}
\label{eq:energy_sideband_cumulant_scaling}
\end{equation}
纯相位 PINEM 的 $P_\ell=J_\ell^2(2|\beta|)$ 关于 $\ell=0$ 对称，所以所有奇阶累积量都严格为零；出现非零偏度意味着入射谱不对称、有限反冲、耗散/非平衡环境或其它打破 $\ell\leftrightarrow-\ell$ 对称的物理机制进入，而不能由理想单频纯相位模型解释。

\subsection{有限作用窗：从振幅积累算到黄金律误差}
\label{sec:electron-finite-window-worked}

上一节的矩形脉冲给出 $\sinc^2$ 跃迁谱。它来自不同时刻发生的跃迁振幅相加：共振时相位相同，失谐时相位随时间转动。为了比较不同脉冲形状，把这一积分保留下来。设初态到末态通道 $j$ 的相互作用绘景矩阵元为 $\hbar\mathcal G_jw(t)\ee^{\ii\Delta_jt}$，其中 $\mathcal G_j$ 是耦合率，$\Delta_j$ 是失谐，两者单位均为 ${\rm s^{-1}}$；无量纲函数 $w(t)$ 描述耦合何时打开、何时关闭以及怎样变化。由一阶含时微扰论，
\begin{equation}
 A_j^{(1)}=-\ii\mathcal G_j\int\dd t\,w(t)\ee^{\ii\Delta_jt},\qquad
 P_j^{(2)}=|\mathcal G_j|^2\left|\int\dd t\,w(t)\ee^{\ii\Delta_jt}\right|^2.
 \label{eq:electron-window-general-amplitude}
\end{equation}
这里已把玻色占据因子或内部矩阵元吸收到所指定通道的 $\mathcal G_j$ 中。上标表示振幅和概率的微扰阶数，而不是交换光子数。它要求总跃迁概率远小于一；不能在某个单通道概率很小时就忽略所有通道累积造成的耗尽。

矩形窗口 $w(t)=1$（$0<t<T$）给出
\begin{equation*}
 A_j^{(1)}=-\ii\mathcal G_jT\ee^{\ii\Delta_jT/2}
 \sinc(\Delta_jT/2).
\end{equation*}
在共振处各时刻同相累积，振幅正比于 $T$，概率正比于 $T^2$；失谐时较晚的振幅逐渐转向，第一次完全相消出现在 $|\Delta_j|=2\pi/T$。概率峰的角频率全宽半高为 $5.566/T$，相应能量全宽半高为 $5.566\hbar/T$。这个宽度属于有限\emph{相干作用窗}，不是衰变态寿命，也不是谱仪分辨率。

\begin{exbox}[高斯作用窗的线宽与等面积比较]
把矩形开关换成 $w(t)=\exp[-t^2/(2\tau_w^2)]$，求弱跃迁谱，并与具有相同积分面积的矩形窗比较。这里 $\tau_w$ 是振幅包络的时间尺度。

\textbf{解。} 对指数配方，或直接使用高斯 Fourier 积分，有
\begin{equation*}
 \int_{-\infty}^{\infty}\dd t\,
 \ee^{-t^2/(2\tau_w^2)+\ii\Delta t}
 =\sqrt{2\pi}\tau_w\ee^{-\Delta^2\tau_w^2/2}.
\end{equation*}
因此
\begin{equation*}
 A_j^{(1)}=-\ii\mathcal G_j\sqrt{2\pi}\tau_w
 \ee^{-\Delta_j^2\tau_w^2/2},\qquad
 P_j^{(2)}=2\pi|\mathcal G_j|^2\tau_w^2
 \ee^{-\Delta_j^2\tau_w^2}.
\end{equation*}
令谱值降到共振峰的一半，取对数得 $|\Delta_j|\tau_w=\sqrt{\ln2}$，故角频率全宽半高为 $2\sqrt{\ln2}/\tau_w$。它来自\emph{概率}包络，不能直接使用振幅包络的半高宽。

若矩形窗长度取 $T=\sqrt{2\pi}\tau_w$，两窗具有相同积分面积和相同共振概率 $|\mathcal G_j|^2T^2$。高斯窗的谱全宽为 $2\sqrt{2\pi\ln2}/T$；矩形窗则为 $4x_{1/2}/T$，其中 $x_{1/2}$ 是 $\sinc^2x=1/2$ 的首个正根。两者都按 $T^{-1}$ 缩窄，但高斯谱没有旁瓣，矩形谱有代数衰减的旁瓣。延长窗口时，弱跃迁条件仍须独立检查，不能只比较线宽。
\end{exbox}

\begin{exbox}[双脉冲的 Ramsey 条纹：延迟如何进入分辨率]
考虑两段等长矩形作用窗，长度为 $\tau_p$，中心分别在 $-D_p/2$ 与 $D_p/2$，其中 $D_p\ge\tau_p$。第二段耦合比第一段多一个可控相位 $\phi_p$，即使用复开关
\begin{equation*}
 w(t)=\mathbf1_{|t+D_p/2|<\tau_p/2}
       +\ee^{\ii\phi_p}\mathbf1_{|t-D_p/2|<\tau_p/2}.
\end{equation*}
符号 $\mathbf1$ 在指定区间内取一、区间外取零。假定两脉冲之间相干保持，并且总跃迁仍可用一阶振幅处理。求总跃迁概率，判断脉冲间隔和单脉冲宽度分别控制哪一部分谱结构。

\textbf{解。} 每段积分都是同一个 $\tau_p\sinc(\Delta\tau_p/2)$，只差中心时刻与参考相位。先相加振幅，得
\begin{align*}
 A^{(1)}&=-\ii\mathcal G\tau_p\sinc(\Delta\tau_p/2)
 [\ee^{-\ii\Delta D_p/2}+\ee^{\ii\phi_p+\ii\Delta D_p/2}]\\
 &=-2\ii\mathcal G\tau_p\sinc(\Delta\tau_p/2)
 \ee^{\ii\phi_p/2}\cos[(\Delta D_p+\phi_p)/2].
\end{align*}
因此
\begin{equation*}
 P^{(2)}=4|\mathcal G|^2\tau_p^2\sinc^2(\Delta\tau_p/2)
 \cos^2[(\Delta D_p+\phi_p)/2].
\end{equation*}
单脉冲时宽决定宽包络，脉冲中心间隔决定包络内的细条纹。$\phi_p=0$ 且 $D_p\gg\tau_p$ 时，中央条纹内的 $\sinc$ 近似为一；由 $\cos^2(\Delta D_p/2)=1/2$ 得全宽半高约为 $\pi/D_p$。因此长延迟可以缩窄条纹，而不增加两段实际作用时间。

若 $\phi_p$ 在重复实验间均匀随机，$\overline{\cos^2}=1/2$，平均概率退化为两个单脉冲概率之和，细条纹消失。把概率先相加相当于预先作了这一失相干假设；它不是相干双脉冲实验的一般规则。
\end{exbox}

单个共振末态的概率随 $T^2$ 增长，而黄金律给出的总概率随 $T$ 增长。要看清这一变化，需要把连续末态在不同失谐处的贡献积分起来。将末态数目与耦合权重合并，定义所选初态的\emph{跃迁谱密度}
\begin{equation}
 \mathcal J_{\rm tr}(\Delta)=\sum_j|\mathcal G_j|^2\delta(\Delta-\Delta_j),\qquad
 P^{(2)}(T)=\int\dd\Delta\,\mathcal J_{\rm tr}(\Delta)
 T^2\sinc^2(\Delta T/2).
 \label{eq:electron-window-transition-spectrum}
\end{equation}
$\mathcal J_{\rm tr}$ 的单位为 $\mathrm{s^{-1}}$，它同时包含末态数目与每个末态的耦合强度。窗口变长时，共振峰高度按 $T^2$ 增长，但参与的连续末态频率宽度按 $1/T$ 缩小，两者相乘才给线性的总概率。对孤立离散态没有这一连续谱积分，不能直接使用恒定黄金律速率。

有限时间误差可以沿用前面洛伦兹谱例题的同一模型。
改记 $\tau_c=a^{-1}$、$J_0=g_*^2/(\pi a)$，则其谱密度写成
\begin{equation*}
 \mathcal J_{\rm tr}(\Delta)=\frac{J_0}{1+\Delta^2\tau_c^2},
 \qquad J_0>0,\qquad \tau_c>0.
\end{equation*}
$J_0$ 是谱峰值，$\tau_c$ 是相关时间。将上述参数转换代入已有二阶结果，得到
\begin{equation}
 P^{(2)}(T)=\Gamma_0[T-\tau_c(1-\ee^{-T/\tau_c})],\qquad
 \Gamma_0=2\pi J_0.
 \label{eq:electron-window-lorentzian-probability}
\end{equation}
短时间给 $P^{(2)}\simeq\Gamma_0T^2/(2\tau_c)$；长时间给 $P^{(2)}\simeq\Gamma_0(T-\tau_c)$。黄金律 $P\simeq\Gamma_0T$ 因而同时需要 $\tau_c/T\ll1$ 和 $\Gamma_0T\ll1$：前者使连续谱有时间去相位，后者使初态尚未明显耗尽。

\begin{exbox}[黄金律的适用时间窗是否存在]
对上述 Lorentz 跃迁谱，给定二阶谱结果相对黄金律的允许误差 $\eta$ 和最大耗尽概率 $p_*$，求一组便于手算的充分条件。

\textbf{解。} 先无量纲化，令 $x=T/\tau_c$、$\epsilon_c=\Gamma_0\tau_c$。精确二阶结果与黄金律之比为
\begin{equation*}
 R(x)=\frac{P^{(2)}(T)}{\Gamma_0T}
 =1-\frac{1-\ee^{-x}}{x},\qquad
 1-R(x)=\frac{1-\ee^{-x}}{x}<\frac1x.
\end{equation*}
这里把相对误差定义为相对于黄金律预测的偏差。求导得
\begin{equation*}
 R'(x)=\frac{1-(1+x)\ee^{-x}}{x^2}>0,
\end{equation*}
因为 $\ee^x>1+x$。因此增加作用时间会单调改善黄金律的有限窗误差，却同时增加耗尽。

由上界，$T\ge\tau_c/\eta$ 足以使二阶谱结果相对黄金律的误差不超过 $\eta$。
耗尽还可以用前面同一连续谱模型的精确振幅检验。
其相关核 $C(s-s')$ 实且非负，联合归一化又给 $|A(s)|\le1$，所以
\[
P_{\rm exact}(T)
=\int_0^T\!\dd s\int_0^T\!\dd s'\,
C(s-s')A(s)A(s')^*
\le P^{(2)}(T)<\Gamma_0T.
\]
因而 $T\le p_*/\Gamma_0$ 确能控制该模型的真实耗尽。
这个概率上界利用了本例相关核非负，不能直接用于任意带振荡的谱。
两条件能够同时满足，只需
\begin{equation*}
 \frac{\tau_c}{\eta}\le T\le\frac{p_*}{\Gamma_0},\qquad
 \Gamma_0\tau_c\le\eta p_*.
\end{equation*}
这是一组充分条件，未必是最宽的时间窗。它直接说明：只有环境去相位时间显著短于初态耗尽时间，才存在可用的恒定速率区。若区间为空，应回到有限时间概率或完整动力学，而不是强行定义一个黄金律速率。
\end{exbox}



\begin{exbox}[同一连续谱中的直接路径与共振路径：Fano 线形]
投影自能确定共振位置和宽度，却还未确定观测线形。
若探测器同时收集直接连续通道与离散共振衰变通道，就必须求两条路径的共同末态振幅。
对下面的单连续通道模型，计算两条路径的相干跃迁线形，并确定反共振零点和峰值。
取能量归一化连续态 $\langle E|E'\rangle=\delta(E-E')$，
\[
H=E_d|d\rangle\langle d|+\int E|E\rangle\langle E|\dd E
+\int[w(E)|E\rangle\langle d|+w(E)^*|d\rangle\langle E|]\dd E.
\]
$w$ 的单位为能量的平方根；态密度已包含在这个归一化中。
令激发算符作用于初态给
$D|i\rangle=\mu_d|d\rangle+\int M(E)|E\rangle\dd E$。
\textbf{解。} 离散预解式为
\begin{align}
G_d^R(E)&=[E-E_d-\Sigma^R(E)]^{-1},
\label{eq:feqo-fano-selfenergy}\\
\Sigma^R(E)&=\mathop{\rm PV}\int\frac{|w(E')|^2}{E-E'}\dd E'
-\ii\pi|w(E)|^2\equiv\Delta_d(E)-\ii\Gamma_d(E)/2,
\label{eq:feqo-fano-width-shift}
\end{align}
其中 $\Gamma_d=2\pi|w|^2$ 为\emph{能量宽度}，衰减率是 $\Gamma_d/\hbar$。

把同一个预解式作用于激发源，先消去连续分量，再投影到出射连续态，得到
\[
T(E)=M(E)+w(E)G_d^R(E)
\left[\mu_d+\int\frac{w(E')^*M(E')}{E-E'+\ii0}\dd E'\right].
\]
这里的积分不能省掉：直接激发连续谱以后，它仍可散射到离散态再返回。
将其分成主值与在壳部分，令
\[
\mu_{\rm eff}(E)=\mu_d+\mathop{\rm PV}\int
\frac{w(E')^*M(E')}{E-E'}\dd E'.
\]
于是括号是 $\mu_{\rm eff}-\ii\pi w^*M$。
在单一连续通道、具有相应时间反演约束且跃迁矩阵元可取实数的模型中，
假定各慢变系数在共振宽度内可取常数，定义
\begin{equation}
\varepsilon=\frac{2(E-E_d-\Delta_d)}{\Gamma_d},
\qquad q_F=\frac{2w\mu_{\rm eff}}{M\Gamma_d}.
\label{eq:feqo-fano-q}
\end{equation}
将两条路径置于同一个分母，得到
\begin{align*}
T(E)&=M+\frac{w\mu_{\rm eff}-\ii\Gamma_dM/2}
{E-E_d-\Delta_d+\ii\Gamma_d/2}\\
&=M\frac{\varepsilon+q_F}{\varepsilon+\ii}.
\end{align*}
因此
\begin{equation}
|T(E)|^2=|M|^2\frac{(\varepsilon+q_F)^2}{1+\varepsilon^2}.
\label{eq:feqo-fano-lineshape}
\end{equation}
分子的虚项被直接通道的再散射精确抵消。
若只加 $M+\mu_dwG_d^R$，就会留下错误的分子常数 $+1$。
$|M|^2$ 是同一耦合连续通道的平滑背景，并非另一个互不干涉的连续谱。
有额外未相干收集的通道或真正损耗时，其强度应另外相加；
复杂相位通常使 $q_F$ 为复数，反共振未必是严格零点。
以上计算是 Fano 构型相互作用的单通道形式。\footnote{U.~Fano 的原始推导见
\url{https://doi.org/10.1103/PhysRev.124.1866}。}

对实数 $q_F\ne0$，求导得
\[
\frac{\dd}{\dd\varepsilon}\frac{(\varepsilon+q_F)^2}{1+\varepsilon^2}
=\frac{2(\varepsilon+q_F)(1-q_F\varepsilon)}{(1+\varepsilon^2)^2}.
\]
零点在 $\varepsilon=-q_F$，峰在 $\varepsilon=1/q_F$，峰高为背景的 $1+q_F^2$ 倍。
例如给定模型参数 $\Gamma_d=\SI{0.15}{\electronvolt}$、$q_F=2$，
零点在共振以下 $\SI{0.15}{\electronvolt}$，峰在共振以上
$\SI{0.0375}{\electronvolt}$，高度为背景五倍。
这组数只是线形模型的代入，不指定某一种材料都满足相同宽度。
\end{exbox}


\subsection{多量子交换与中间态预解式}

上一小节的一阶振幅只包含一次量子交换。若作用更强或时间更长，同一个电子还可以先交换一个量子，再交换第二个量子。含时演化的 Dyson 级数已经包含这些过程：每插入一次相互作用，就有一次产生或湮灭算符作用；二阶项包含两次作用，并对先后时刻积分。本节求出这些中间态的贡献，为下一章的多边带分布作准备。

先保留两个任意环境模式 $\lambda,\mu$，频率分别为 $\omega_\lambda,\omega_\mu$，电子中间态仍取完整连续谱。相互作用绘景的线性耦合写成
\[
\hat V_I(t)=\sum_\nu[\hat V_{\nu,+}(t)+\hat V_{\nu,-}(t)],
\]
其中 $\hat V_{\nu,+}$ 使模式 $\nu$ 的占据数增加一，$\hat V_{\nu,-}=\hat V_{\nu,+}^\dagger$。若 $\lambda\neq\mu$，从 $|i,n_\lambda,n_\mu\rangle$ 到 $|f,n_\lambda+1,n_\mu+1\rangle$ 的最低非零振幅是二阶，并包含两种时间排序：
\begin{equation}
{
\begin{aligned}
\mathcal A_{fi;\lambda\mu}^{(2,\,2{\rm em})}
={}&\left(-\frac{\ii}{\hbar}\right)^2
\int_{t_i}^{t_f}\dd t_2\int_{t_i}^{t_2}\dd t_1\sum_m
\Big[
\langle f|\hat V_{\lambda,+}(t_2)|m\rangle
\langle m|\hat V_{\mu,+}(t_1)|i\rangle\\
&\hspace{34mm}+
\langle f|\hat V_{\mu,+}(t_2)|m\rangle
\langle m|\hat V_{\lambda,+}(t_1)|i\rangle
\Big].
\end{aligned}}
\label{eq:general_sequential_two_boson_dyson_amplitude}
\end{equation}
这里为简洁省略了外部 Fock 标签；玻色矩阵元整体给出
$\sqrt{(n_\lambda+1)(n_\mu+1)}$。指标 $m$ 表示完整电子中间谱，对自由电子就是自旋求和与连续动量积分，而不是预先挑选的有限维“虚能级”。

在绝热开关与长时间散射极限，两种时间排序分别给出不同的中间态能量分母：
\begin{equation}
{
\begin{aligned}
\mathcal M_{fi;\lambda\mu}^{(2,\,2{\rm em})}
={}&\sqrt{(n_\lambda+1)(n_\mu+1)}\sum_m
\left[
\frac{V^{(\lambda)}_{fm}V^{(\mu)}_{mi}}
{E_i-E_m-\hbar\omega_\mu+\ii0^+}
+
\frac{V^{(\mu)}_{fm}V^{(\lambda)}_{mi}}
{E_i-E_m-\hbar\omega_\lambda+\ii0^+}
\right],\\
&\hspace{12mm}E_f-E_i+\hbar(\omega_\lambda+\omega_\mu)=0.
\end{aligned}}
\label{eq:two_boson_resolvent_amplitude_general}
\end{equation}
若电子中间谱连续，$\sum_m$ 应理解为 $\sum_s\int\dd^3p$。每个分母的 $\ii0^+$ 同时规定主值虚跃迁与真实在壳中间态的解析边界；一旦某个中间态落入有限时间线宽，就应把它提升为显式动力学通道，而不是继续只保留为离壳预解式分母。

两个量子属于同一个模式时，令 $\lambda=\mu$ 并合并相同的时间排序即可。外部玻色因子变成
$\sqrt{(n+1)(n+2)}$，总能量守恒退化为
$E_f-E_i+2\hbar\omega_0=0$。因此常见的单模双光子过程只是上式的退化情形，而不是多量子交换的母定义。

更一般地，$N$ 个线性顶点产生的净多量子振幅是所有时间排序、中间电子路径和发射/吸收序列的相干和。因而“$N$ 光子过程”首先是散射路径的阶数与净量子数标签，而不是自动意味着哈密顿量中存在一个 $a^N$ 基本顶点。只有当母理论本身含有二次正常坐标耦合，或者把高能/负能中间态作 Feshbach--Schur 消元后生成了局域二次有效算符时，才会出现真正的直接二量子顶点
\begin{equation}
{
\hat V_{\rm dir}^{(2)}
=\frac{\hbar}{2}
\sum_{\lambda\mu}
\hat{\mathcal Q}_{\lambda\mu}
(\hat a_\lambda+\hat a_\lambda^\dagger)
(\hat a_\mu+\hat a_\mu^\dagger),}
\label{eq:general_direct_quadratic_boson_vertex}
\end{equation}
其中 $\hat{\mathcal Q}_{\lambda\mu}$ 是电子空间上的频率算符。此时同一个两量子末态既收到\eqnrefs{eq:general_sequential_two_boson_dyson_amplitude} 的顺序二阶贡献，也收到\eqnrefs{eq:general_direct_quadratic_boson_vertex} 的一阶直接贡献，两者必须在振幅层相加后再取模平方。

电磁相互作用还需要区分“基本顶角”和“投影后的有效顶角”。完整 Dirac 最小耦合的 QED 基本顶角线性于 $A_\mu$。正能低能投影或 Foldy--Wouthuysen 约化后出现的 $A^2$ 项，已经包含被消去相对论中间态的影响，因此不能再把它当作与两个线性 QED 顶角无关的第二种基本作用。

Kapitza--Dirac 的有质动力势、两光子 Bragg 路径和完整 QED 中的顺序双顶角振幅可以描述同一低能物理，但它们属于不同约化层级；计算时必须选定一个层级，避免把同一中间态贡献重复计数。

\begin{derivation}[有限作用时间的一量子发射概率]
对矩形作用窗完成时间积分并取模平方，就得到正文\eqnrefs{eq:universal_boson_finite_time_emission} 的 $\sinc^2$ 谱核。 在相互作用绘景中，电子算符与模式算符分别按自由哈密顿量 演化。 对电子本征态，

\begin{align*}
{}_V\!\langle f_e|\delta\hat h_\lambda^\dagger(t)|i_e\rangle_V
&=\ee^{\ii(E_f-E_i)t/\hbar}
{}_V\!\langle f_e|\delta\hat h_\lambda^\dagger|i_e\rangle_V,\\
\langle n_\lambda+1|\hat a_\lambda^\dagger(t)|n_\lambda\rangle
&=\ee^{+\ii\omega_\lambda t}\sqrt{n_\lambda+1}.
\end{align*}
因此发射末态 $|f_e;\mathbf n+\mathbf e_\lambda\rangle_V$ 与初态之间的瞬时矩阵元为
\begin{align*}
\langle f_e;\mathbf n+\mathbf e_\lambda|\hat H_I(t)|i_e;\mathbf n\rangle
&=M_{fi,\lambda}^{\rm em}\sqrt{n_\lambda+1}
\exp\!\left\{\ii\left[\frac{E_f-E_i}{\hbar}+\omega_\lambda\right]t\right\}\\
&=M_{fi,\lambda}^{\rm em}\sqrt{n_\lambda+1}
\ee^{\ii\Delta_{fi}^{\rm em}t}.
\end{align*}
Dyson 级数的一阶项于是给
\[
C_f^{(1)}(T)
=-\frac{\ii}{\hbar}M_{fi,\lambda}^{\rm em}\sqrt{n_\lambda+1}
\int_0^T\dd t\,\ee^{\ii\Delta_{fi}^{\rm em}t}.
\]
时间积分可以写成
\[
\int_0^T\dd t\,\ee^{\ii\Delta t}
=T\ee^{\ii\Delta T/2}\sinc(\Delta T/2),
\]
故
\[
C_f^{(1)}(T)
=-\ii\mathcal G_{fi,\lambda}^{\rm em}\sqrt{n_\lambda+1}
T\ee^{\ii\Delta_{fi}^{\rm em}T/2}
\sinc(\Delta_{fi}^{\rm em}T/2).
\]
取模平方时中间相位 $\exp(\ii\Delta T/2)$ 消失，得到\eqnrefs{eq:universal_boson_finite_time_emission}。其中 $n_\lambda+1$ 直接来自 $|\langle n_\lambda+1|a_\lambda^\dagger|n_\lambda\rangle|^2$；即使 $n_\lambda=0$，产生算符仍有非零矩阵元。吸收支路的末态为 $|f_e;\mathbf n-\mathbf e_\lambda\rangle_V$。此时
\begin{align*}
{}_V\!\langle f_e|\delta\hat h_\lambda(t)|i_e\rangle_V
&=\ee^{\ii(E_f-E_i)t/\hbar}M_{fi,\lambda}^{\rm abs},\\
\langle n_\lambda-1|\hat a_\lambda(t)|n_\lambda\rangle
&=\ee^{-\ii\omega_\lambda t}\sqrt{n_\lambda},
\end{align*}
所以总相位的角频率失谐为
\[
 \Delta_{fi}^{\rm abs}
 =\frac{E_f-E_i}{\hbar}-\omega_\lambda
 =\frac{E_f-E_i-\hbar\omega_\lambda}{\hbar}.
\]
因此
\begin{align*}
 C_{f,\rm abs}^{(1)}(T)
 &=-\ii\mathcal G_{fi,\lambda}^{\rm abs}\sqrt{n_\lambda}
 T\,\ee^{\ii\Delta_{fi}^{\rm abs}T/2}
 \sinc(\Delta_{fi}^{\rm abs}T/2),\\
 P_{i\to f}^{\rm abs}(T)
 &=n_\lambda|\mathcal G_{fi,\lambda}^{\rm abs}|^2T^2
 \sinc^2(\Delta_{fi}^{\rm abs}T/2).
\end{align*}
这与发射支路的 $n_\lambda+1$ 和 $\Delta_{fi}^{\rm em}$ 构成两条独立的有限时间结果。
\end{derivation}

\begin{derivation}[两次量子发射的时间排序与中间态]
取时间无关耦合在有限散射窗内绝热开关，并写发射顶角的相互作用绘景矩阵元
\[
\langle m,n+1|\hat V_+(t_1)|i,n\rangle
=\sqrt{n+1}\,V^{\rm em}_{mi}
\exp\!\left[
\frac{\ii(E_m-E_i+\hbar\omega_0)t_1}{\hbar}
\right],
\]
\[
\langle f,n+2|\hat V_+(t_2)|m,n+1\rangle
=\sqrt{n+2}\,V^{\rm em}_{fm}
\exp\!\left[
\frac{\ii(E_f-E_m+\hbar\omega_0)t_2}{\hbar}
\right].
\]
代入二阶 Dyson 振幅，定义
\[
\Omega_1=\frac{E_m-E_i+\hbar\omega_0}{\hbar},
\qquad
\Omega_2=\frac{E_f-E_m+\hbar\omega_0}{\hbar},
\]
则时间核为
\[
I_T(\Omega_2,\Omega_1)
=\int\dd t_2\int^{t_2}\dd t_1\,
\ee^{\ii\Omega_2t_2}\ee^{\ii\Omega_1t_1}.
\]
时间排序区域 $t_1<t_2$ 最适合改用“总时间”和“两次顶角的时间间隔”作坐标：
\begin{equation*}
 t\equiv t_2,
 \qquad
 \tau\equiv t_2-t_1\ge0,
 \qquad
 t_1=t-\tau,
 \qquad
 \dd t_2\,\dd t_1=\dd t\,\dd\tau.
\end{equation*}
指数因子于是精确分解成
\begin{align*}
 \ee^{\ii\Omega_2t_2}\ee^{\ii\Omega_1t_1}
 &=\ee^{\ii\Omega_2t}
   \ee^{\ii\Omega_1(t-\tau)}\\
 &=\ee^{\ii(\Omega_1+\Omega_2)t}
   \ee^{-\ii\Omega_1\tau}.
\end{align*}
$t$ 积分只检查整个过程的总能量守恒，而 $\tau\ge0$ 记录第一顶角之后中间态能够传播的持续时间；两个积分因此分别编码整体守恒律与中间态传播。

为把半轴积分定义成分布，在相对时间上加入收敛因子 $\ee^{-\epsilon\tau}$，$\epsilon>0$，最后取 $\epsilon\to0^+$：
\begin{align*}
 I_\epsilon(\Omega_2,\Omega_1)
 &=\int_{-\infty}^{\infty}\dd t\,
   \ee^{\ii(\Omega_1+\Omega_2)t}
   \int_0^\infty\dd\tau\,
   \ee^{-(\epsilon+\ii\Omega_1)\tau}\\
 &=2\pi\delta(\Omega_1+\Omega_2)
   \frac1{\epsilon+\ii\Omega_1}\\
 &=2\pi\delta(\Omega_1+\Omega_2)
   \frac{-\ii}{\Omega_1-\ii\epsilon}.
\end{align*}
总时间部分给
\begin{equation*}
2\pi\delta(\Omega_1+\Omega_2)
=2\pi\hbar\,
\delta(E_f-E_i+2\hbar\omega_0).
\end{equation*}
另一方面，二阶 Dyson 前因子是 $(-\ii/\hbar)^2=-1/\hbar^2$。因此相对时间积分与 Dyson 前因子的乘积为
\begin{align*}
 -\frac1{\hbar^2}
 \frac{-\ii}{\Omega_1-\ii0^+}
 &=\frac{\ii}{\hbar^2}
 \frac{\hbar}{E_m-E_i+\hbar\omega_0-\ii0^+}\\
 &=-\frac{\ii}{\hbar}
 \frac1{E_i-E_m-\hbar\omega_0+\ii0^+}.
\end{align*}
于是完整二阶 $S$ 矩阵元具有标准的形式
\begin{align*}
 \langle f,n+2|S^{(2)}|i,n\rangle
 ={}&-2\pi\ii\,
 \delta(E_f-E_i+2\hbar\omega_0)\\
 &\times\sqrt{(n+1)(n+2)}
 \sum_m
 \frac{V^{\rm em}_{fm}V^{\rm em}_{mi}}
 {E_i-E_m-\hbar\omega_0+\ii0^+}.
\end{align*}
去掉散射矩阵中共同的 $-2\pi\ii\delta(E_f-E_i+2\hbar\omega_0)$ 后，约化二阶振幅正是
\[
\mathcal M_{fi}^{(2,\,2{\rm em})}
=\sqrt{(n+1)(n+2)}
\sum_m
\frac{V^{\rm em}_{fm}V^{\rm em}_{mi}}
{E_i-E_m-\hbar\omega_0+\ii0^+}.
\]
若电子中间谱连续，$\sum_m$ 展开为对完整动量、自旋和其它量子数的积分。利用
\[
\frac1{x+\ii0^+}
=\mathcal P\frac1x-\ii\pi\delta(x),
\]
可见中间分母的主值部分对应离壳虚路径，而 $\delta(x)$ 部分表示中间态本身可以在壳传播。后一情形若在有限脉冲带宽内占显著权重，就应恢复显式级联通道，而不应继续把全过程压成局域二阶有效顶角。
\end{derivation}

\begin{derivation}[有限时间谱核与 Lorentz 线形的积分]
保留有限时间的谱核，对 Lorentz 线形直接积分，以求出概率从短时间到长时间的变化。将谱核写成两时间积分，$T^2\sinc^2(\Delta T/2)=\int_0^T\dd t\int_0^T\dd t'\,\ee^{\ii\Delta(t-t')}$，再使用 Lorentz 线形的 Fourier 积分
\begin{equation*}
 \int_{-\infty}^{\infty}\dd\Delta\,
 \frac{J_0\ee^{\ii\Delta u}}{1+\Delta^2\tau_c^2}
 =\frac{\pi J_0}{\tau_c}\ee^{-|u|/\tau_c}.
\end{equation*}
两时间正方形区域中，相同时间差 $u>0$ 对应的线段长度为 $T-u$，故
\begin{align*}
 P^{(2)}(T)&=\frac{2\pi J_0}{\tau_c}
 \int_0^T(T-u)\ee^{-u/\tau_c}\dd u\\
 &=2\pi J_0[T-\tau_c(1-\ee^{-T/\tau_c})].
\end{align*}
因此恒定速率是有限时间积分的渐近结果。谱阈值、窄峰或离散末态改变了被积分的 $\mathcal J_{\rm tr}$ 时，应先重算这个积分，再判断是否存在黄金律区。
\end{derivation}

